\documentclass{article}
\usepackage{graphicx} % Required for inserting images
\usepackage{ctex}

\usepackage{amsfonts}
\usepackage{amssymb}
\usepackage{amsmath}
\usepackage{amsthm}
% 导入 xparse 宏包以支持 LaTeX3 语法
\usepackage{xparse}
\usepackage{pgfplots}
% 用于插入带有坐标轴、标签和曲线的图像

% \begin{tikzpicture}
%     \begin{axis}[
%       xlabel=$x$,
%       ylabel=$f(x)$,
%       axis lines=middle,
%       xmin=-5, xmax=5,
%       ymin=-2, ymax=8,
%       width=\textwidth,
%       height=8cm
%     ]
%     \addplot[blue,domain=-3:3] {x^2};
%     \end{axis}
%   \end{tikzpicture}
\usepackage{tikz}
% 用于绘制一般的图像

% \begin{tikzpicture}[scale=0.8]
%     \draw[->] (-4,0) -- (4,0) node[right] {$x$};
%     \draw[->] (0,-1) -- (0,9) node[above] {$f(x)$};
%     \draw[domain=-3:3,smooth,variable=\x,blue] plot ({\x},{\x^2});
%   \end{tikzpicture}

\newtheorem{theorem}{定理}[subsection]
\newtheorem{lemma}{引理}[subsection]
\newtheorem{corollary}{推论}[subsection]
\newtheorem{example}{例}[subsection]

% 为了证明中可以使用中文，后续定义证明时使用cproof而不是proof
\newenvironment{cproof}{%
\heiti{证明}\kaishu
}{%
%   \hfill $\square$ 添加结束符号
%   \par\bigskip 可选的垂直间距
}
\newenvironment{identification}{%
\heiti{定义}\kaishu
}{%
%   \hfill $\square$ 添加结束符号
%   \par\bigskip 可选的垂直间距
}


\newcommand{\RR}{\mathbb{R}}
\newcommand{\NN}{\mathbb{N}}
\newcommand{\CC}{\mathbb{C}}
\newcommand{\QQ}{\mathbb{Q}}
\newcommand{\ZZ}{\mathbb{Z}}
\newcommand{\FF}{\mathbb{F}}
% 简化各种常见数的集合

\newcommand{\parameter}[1]{\left(#1\right)}

\newcommand{\bracket}[1]{\left[#1\right]}

\newcommand{\abs}[1]{\left|#1\right|}

\newcommand{\mo}[1]{\|#1\|}
% 各种自动变化大小的括号的简化

\newcommand{\ve}{\boldsymbol}
% 为了适应David C Lay线性代数中，简化斜体+粗体向量的书写

\newcommand{\base}{\mathcal}

\newcommand{\tb}{\textbf}
% 简化粗体字体的书写

\newcommand{\f}[2]{\frac}
\newcommand{\df}[2]{\dfrac}

\NewDocumentCommand{\vs}{m m m}{
    \boldsymbol{#1}_#2,\cdots,\boldsymbol{#1}_#3
}
% 快速书写一个向量组，第一个参数为向量名称，后两个为首末角标

% $\vs{b}{1}{n} $这是多参数命令的使用示例

\NewDocumentCommand{\cvs}{m m m m m}{
    #1_{#4}\boldsymbol{#2}_#4 #3 \cdots #3 #1_{#5}\boldsymbol{#2}_#5
}
% 快速书写一个向量线性组合，第一个参数为系数，第二个参数为向量名称，第三个参数为运算符，后两个参数为角标

\title{数学分析}
\author{徐海翁}
\date{2023.9.11}

\begin{document}
\begin{CJK}{UTF8}{gkai}
\maketitle

\subsection{课程概论与学习方法}
习题课会留下思考题，可以思考一个星期\\
作业本2本轮换使用\\
作业和上课基本与数院等同\\
考试难度不难于平时作业\\
数分：内容很多，第一时间领悟到\\
解决问题的习惯：问题不过夜，当天讲的问题，当天解决（问老师、助教、同学）\\
做作业的习惯：完全弄清楚本课的内容后再开始，完成作业时不要翻书、翻笔记\\
现代数学的基础：集合论\\
进而形成的几大分支：\\
元素的运算——代数学；\\
元素之间的关系（位置关系）——拓扑学（几何学）；\\
集合之间的映射——分析学（最大的一个分支，起点不高，最广泛被接受）\\
\newpage
\section{实数}
\subsection{集合论知识补充}
\subsubsection{集合的表示方法}
\paragraph{列举法\\}
\paragraph{描述法\\}
形如$A=\{x|x+3=1\}$\\
\subsection{数学分析的研究对象}
函数及其性质（分析性质）\\
什么是函数?一种特殊的映射（$X\to Y$的映射，其中$X$，$Y$均为实数（复数），则为实（复变）函数，$X$为实数集，$Y$为复数集——复值函数\\
函数的内涵：定义域+对应关系$f$【$\sin$是函数而$\sin x$不是函数，后者只是一个函数值】，其他东西如$x$，$y$等都只属于记号！\\
近代数学与古代数学的最大差异——函数的出现（几乎出现在数学的任何一个领域）\\
古代使用了具体的函数，但没有把函数的抽象概念提出\\
数学分析不研究具体函数的性质，而是研究所有函数的共同性质（研究的是很宏大的东西）\\
\subsubsection{函数的表示方法}
\paragraph{图表\\}
\paragraph{曲线\\}
\paragraph{表达式（使用很多）\\}
如果$f$和$g$的定义域相同（均记为$A$），且他们的对应关系$A\rightarrow B$也相同，则$f$与$g$为同一函数（不用写成$f(x)$这种）
\paragraph{有尽小数在实数集上是稠密的\\}
对于任意两个规范小数$a$,$b$
\[a=a_0. a_1 a_2 a_3\dots\]
\[b=a_0. a_1 a_2 a_3\dots\]
\[a<b,\exists p \in {N} s.t. a_p<b_p\]
\[\because a,b\text{是规范小数} \]
\[\therefore  \exists q>p \quad  s.t.\quad a_q<9\]
\[\therefore \exists c=a_0. a_1 a_2 \dots (a_q+1) s.t. a<c<b\]
\subsubsection{确界原理}
$E=\{p\in \mathbb{Q}  |p^2<2\}$有界但没有“边界”\\

$E=\{p\in \mathbb{R}  |p^2<2\}$有界且有上确界\\
\subsection{可数集和不可数集}
\paragraph{等势：称集合$A,B$等势，若存在$A$，$B$之间的一一映射。\\}
\paragraph{可数集：与自然数集等势的集合称为可数集（元素可被自然数标号）\\}
举例：偶数集与自然数集等势\\
因为可以建立一一对应的关系第$k$号元素为$2(k+1)$\\
\paragraph{两个可数集的并集亦为可数集\\}
\subsubsection{有理数是可数集}
\[x_1^1 \quad  x_2^1 \quad  x_3^1 \quad  \dots\]
\[x_1^2 \quad  x_2^2 \quad  x_3^2 \quad  \dots\]
\[x_1^3 \quad  x_2^3 \quad  x_3^3 \quad  \dots\]
第一行为分母为1的分数：$$1 \quad \frac{2}{1} \quad \frac{3}{1} \quad \dots$$ \newline
第二行为分母为2的分数：$$\frac{1}{2} \quad \frac{2}{2} \quad \frac{3}{2} \quad \dots$$ \newline
可以被一个个数完：从左上角到右下角\\
\subsubsection{实数集是不可数集（反证法）}
假设我能够遍历所有实数给他们编号，那么可以构造一个新的实数，这个实数的第i位都与已有的第i个实数的第i位不同，那么显然这个实数不在已经遍历的实数内，然而这与我们已经遍历的所有的实数矛盾，故假设不成立\\
if we can list all the real numbers,
for example , at any time ,for the first $r$ numbers ,
\[a_1=a_1^1.a_1^2  a_1^3  a_1^4  \dots\]
\[a_2=a_2^1.a_2^2  a_2^3  a_2^4  \dots\]
\[a_3=a_3^1.a_3^2  a_3^3  a_3^4  \dots\]
\[\dots\]

\[a_r=a_r^1.   a_r^2   a_r^3   a_r^4  \dots\]
we could construct a new number $b$,which has the characteristics below
\[(1) \quad  b=b_1.b_2b_3b_4\dots\]
\[(2) \quad  \forall i \in \mathbb{N} \quad  b_i\neq a_i^i \dots\]
then $b$ is unequal to all the real numbers we had already listed, but what we had supposed is that we had already listed all the real numbers.That's an obvious contradiction!
So the assumption is not correct.

\newpage
\section{极限}
\subsection{有界序列与无穷小序列}
\subsubsection{有界序列}
\paragraph{定义\\}

设$\{x_n\}$是一个序列\\
\[\exists M\in\mathbb{R},s.t x_n\leq M ,\forall n\in \mathbb{N}\]
则称$\{x_n\}$有上界，实数$M$是它的一个上界
\[\exists m\in\mathbb{R},s.t x_n\geq m m\forall n\in \mathbb{N}\]
则称$\{x_n\}$有下界，实数$m$是它的一个上界\\
若同时有上界与下界，则称序列有界\\
序列有界的表述：
\[\exists K\in \mathbb{R},\forall n\in \mathbb{N},s.t.\quad x_n\leq |K|\]
相应的，序列无界的表述
\[\forall K\in \mathbb{R},\exists n\in \mathbb{N},s.t.\quad  x_n  >  |K|\]

这里我们可以对比“无界”与“收敛于无穷”的区别，前者是一个序列的整体性质，而后者表示一个趋近的过程\\
无界的反面是有界，但发散于无穷的反面并非收敛于有限值（这在反证时一定要区分开来）\\
\paragraph{例子\\}

$x_n=n\sin n$是有界序列吗？
\[\forall N \in \mathbb{N} ,\exists n> \mathbb{N} s. t. n \in (2k\pi+\frac{\pi}{3},2k\pi+\frac{2\pi}{3})\quad  s.t. \quad  sin n>\frac{\sqrt{3}}{2}\]
\[\therefore \exists N_0=2M \quad  s.t.\quad \exists n>N_0,nsin n > \sqrt{3}M>M\]
\[\therefore \text{序列} \{x_n\} \text{是无界序列}\]


\subsubsection{无穷小序列}
\paragraph{定义\\}
设$\{x_n\}$是一个实数序列，对$\forall \varepsilon > 0  $,$\exists N\in \mathbb{N}$使得$\forall n > N$有$|x_n|<\varepsilon$,则该序列为\textbf{无穷小序列}\\
用符号写作
\[\forall \varepsilon > 0 ,\exists N\in \mathbb{N},\forall n > N, |x_n|  <   \varepsilon \]
相应的，序列是无穷大序列可以表示为\\
\[\exists \varepsilon > 0 ,\forall N\in \mathbb{N},\exists n > N, |x_n| \geq \varepsilon \]
\begin{lemma}
设$\{\alpha_n\}$和$\{\beta_n\}$是实数序列，若$\exists N_0\in \mathbb{N}$使得
\[|\alpha_n| < \beta_n,\forall n > N_0\]
如果$\{\beta _n\}$是无穷小序列，则$\{\alpha _n\}$也是无穷小序列\\
\end{lemma}

\subsubsection{有界序列与无穷小序列的性质}

\begin{lemma}
如果$\{\alpha_n\}$是无穷小序列，则它是有界序列\\
\end{lemma}
\paragraph{主要性质\\}
(1)两个有界序列$\{x_n\},\{y_n\}$的和$\{x_n+y_n\}$与积$\{x_n y_n\}$也是有界序列\\
(2)两个无穷小序列$\{\alpha_n\}\{\beta_n\}$的和$\{\alpha_n+\beta_n\}$也是无穷小序列\\
(3)无穷小序列$\{\alpha_n\}$与有界序列$\{x_n\}$的乘积$\{\alpha_n x_n\}$是有界序列\\
(4)$\{a_n\}$是无穷小序列$\Leftrightarrow\{|a_n|\}$是无穷小序列\\
(5)两个无穷小序列$\{\alpha_n\}\{\beta_n\}$的乘积$\{\alpha_n\beta_n\}$也是无穷小序列\\
(6)实数$c$与无穷小序列的乘积$\{c\alpha_n\}$也是无穷小序列\\
(7)有限个无穷小序列的和与乘积也是无穷小序列\\


为什么我们在上面要进行这么多对“有限”的限定？\\
\paragraph{无穷可数个无穷小序列之和}
for example\[\{x_n^1\},\{x_n^2\},\{x_n^3\}\quad \{x_n\}:x_n=x_n^1+x_n^2+x_n^3\dots\]
(1) 没有意义（意思是找不到这样的实数与结果进行对应） 如\[x_n^1=x_n^2=\dots=\frac{1}{n}\]
(2) 仍是无穷小序列\[x_n^k=\frac{1}{2^k}\]
(3) 有界非无穷小序列$x_n^1$的每一项依次为$1\quad 0 \dots 0 0 0$
即$x_n^k$总是第$k$项为$1$，其余的都是$0$，那么所得到的结果是$x_n$的每一项都是$1$，有界但并非无穷小

\paragraph{无穷可数个无穷小序列之积}
for example\[\{x_n^1\},\{x_n^2\},\{x_n^3\}\quad \{x_n\}:x_n=x_n^1·x_n^2·x_n^3\dots\]
(1) 没有意义 如\[\{x_n^1\}:1 \quad\frac{1}{2}\quad \frac{1}{3} \dots\]
\[\{x_n^2\}:2\quad1\quad \frac{1}{2}\quad \frac{1}{3}\dots \]
\[ \{x_n^3\}: 3 \quad2\quad 1\quad \frac{1}{2}\quad \frac{1}{3} \dots \]
(2) 仍是无穷小序列\[x_n^k=\frac{1}{n}\]
(3) 有界非无穷小序列如\[\{x_n^1\}:1 \quad \frac{1}{n}\quad \frac{1}{n} \dots\]
\[\{x_n^2\}:1 \quad n\quad \frac{1}{n}\quad \frac{1}{n}\dots \]
\[\{x_n^3\}:1 \quad 1\quad n^2 \quad \frac{1}{n}\quad \frac{1}{n} \dots \]

\subsection{收敛序列}
\subsubsection{收敛序列的定义}
\paragraph{定义\\}
设$\{x_n\}$是实数序列，$a$是实数，若对于任意实数$\varepsilon > 0$都存在自然数$N$，使得只要$n > N$，就有
\[|x_n-a|<\varepsilon\]
那么我们就说序列$\{x_n\}$收敛于$a$\\

用符号表示为
\[\forall \varepsilon>0,\exists N\in \mathbb{N},\forall n > N, | x_n -a |   <  \varepsilon \]
相应的，序列不收敛于$a$的符号表示为
\[\exists \varepsilon>0,\forall N\in \mathbb{N},\exists n > N, | x_n -a | \geq \varepsilon \]

\begin{theorem}
若序列$\{x_n\}$的极限存在，则该极限是唯一的\\
\end{theorem}

\begin{theorem}（夹逼原理）

设$\{x_n\},\{y_n\},\{z_n\}$都是实数序列，满足条件
\[x_n\leq y_n\leq z_n\quad, \forall n\in \mathbb{N}(\text{in fact,} \exists N\in \mathbb{N},s.t. \forall n >N)\]
如果\[\lim x_n=\lim z_n=a\]
则$\{y_n\}$也是收敛序列，且其收敛于$a$，即
\[\lim_{} y_n = a\]
\end{theorem}

\begin{theorem}
设$\{x_n\}$是实数序列，$a$是实数，则下列三陈述等价
(1)序列$\{x_n\}$以$a$为极限\\
(2)序列$\{x_n-a\}$为无穷小序列\\
(3)存在无穷小序列$\{a_n\}$，使得
\[x_n=a_n+a,n=1,2,\cdots\]
\end{theorem}
\subsubsection{收敛序列的性质}
\begin{theorem}
收敛序列$\{x_n\}$是有界的\\
\end{theorem}
\begin{theorem}
(1)$\lim |x_n|=|\lim x_n|$\\
(2)$\lim(x_n\pm y_n)=\lim x_n \pm \lim y_n$\\
(3)$\lim(x_n y_n)=\lim x_n  \lim y_n$\\
(4)$\lim\dfrac{1}{x_n}=\dfrac{1}{\lim x_n}$\\
(5)$\lim(c x_n)=c\lim x_n$\\
(6)$\lim\dfrac{y_n}{x_n}=\dfrac{\lim y_n}{\lim x_n}$\\
\end{theorem}
\subsubsection{收敛序列与不等式}
\begin{theorem}

如果$\lim x_n <\lim y_n$,那么存在$N\in \mathbb{N}$，使得$n > N$时，有\\
\[x_n < y_n\]
\end{theorem}

\begin{theorem}
如果$\{x_n\},\{y_n\}$都是收敛序列，且满足条件
\[x_n \leq y_n ,\forall n > N_0\]
则有
\[\lim x_n \leq \lim y_n\]
\end{theorem}
\subsection{收敛原理}
\subsubsection{单调收敛原理}
\begin{theorem}
递增序列$\{x_n\}$收敛的充要条件是它有上界\\

递减序列$\{y_n\}$收敛的充要条件是它有下界\\
\end{theorem}
注：这个定理中的递增（递减）条件可以削弱为$(\exists N\in \mathbb{N},\forall n > N)$
不过这个情况下极限未必是整个序列的上（下）确界\\
\subsubsection{闭区间套原理与B-W定理}
如果一列闭区间$[a_n,b_n]$满足如下条件\\
(1)$[ a_{n+1},b_{n+1}]\subset[ a_n,b_n]$\\
(2)$\lim(b_n-a_n)=0$\\
那么我们就称该列闭区间形成\textbf{闭区间套}\\
\begin{theorem}（闭区间套定理）\\

如果序列$\{a_n\},\{b_n\}$满足以下条件\\

(1)$a_n\leq a_{n+1}\leq b_{n+1}\leq b_n$\\

(2)$\lim(b_n-a_n)=0$\\

那么\\

(1)序列$\{a_n\}$与$\{b_n\}$收敛于相同的极限值
\[\lim a_n=\lim b_n=c\]
(2)$c$是满足下面条件的唯一实数
\[a_n\leq c\leq b_n,\forall n\in \mathbb{N}\]
\end{theorem}
\begin{theorem}
设序列$\{x_n\}$收敛于$a$，则它的任何子序列$\{x_{n_k}\}$都收敛于相同的极限值\\
\end{theorem}
\begin{theorem}（B-W定理）\\
如果一个序列$\{x_n\}$有界，那它一定具有收敛的子序列\\
\subsubsection{柯西收敛原理}
\end{theorem}
\paragraph{定义\\}
如果序列$\{x_n\}$满足条件使得$\forall \varepsilon>0,\exists N\in \mathbb{N},\forall m,n>N,|x_m-x_n|<\varepsilon$,我们就称这个序列为柯西序列（基本序列）\\

当然，在一些题目中使用反证时，也可以利用柯西收敛原理的反面
\[\exists \varepsilon>0,\forall N>0,\exists n,m>N,s.t. |x_m-x_n|\geq \varepsilon\]
\begin{theorem}
序列$\{x_n\}$收敛的充分必要条件是
\[\forall \varepsilon>0,\exists N\in \mathbb{N},\forall m,n>N,|x_m-x_n|<\varepsilon\]

这说明$\{x_n\}$是收敛序列$\Leftrightarrow\{x_n\}$是柯西序列\\
\end{theorem}
\subsection{无穷大}

\subsubsection{无穷极限}
\paragraph{定义\\}
这里我们只给一个简单的例子，其他的定义同理\\

对于序列$\{x_n\}$\\
对于任意的实数$E$，总存在一个$N$，使得任意的$n>N$，都有$x_n > E$
我们就称$\{x_n\}$发散于正无穷，记为
\[\lim_{n\to +\infty}x_n=+\infty\]

以下是几个由实数推广到扩充实数系的定理\\
\begin{theorem}(单调收敛原理的推广)\\

若序列$\{x_n\}$单调递增(递减)，则序列广义收敛\\
\end{theorem}
\begin{theorem}(夹逼定理的推广)\\

若序列$\{x_n\},\{y_n\}$且满足$\lim_{n\to +\infty}x_n=+\infty$，且满足$\exists N\in \mathbb{N},\forall n > N$,
\[x_n\leq y_n\]
那么有
\[\lim_{n\to +\infty}y_n=+\infty\]

若序列$\{x_n\},\{y_n\}$且满足$\lim_{n\to +\infty}y_n=-\infty$，且满足$\exists N\in \mathbb{N},\forall n > N$,
\[x_n<y_n\]
那么有
\[\lim_{n\to +\infty}x_n=-\infty\]
\end{theorem}
\begin{theorem}

如果对于序列$\{x_n\}$有
\[\lim_{n\to +\infty}x_n=+\infty\]
则对于$\{x_n\}$任意子序列$\{x_{n_k}\}$也有
\[\lim_{k\to +\infty}x_{n_k}=+\infty\]
\end{theorem}
\begin{theorem}

$\{x_n\}$是无穷小序列$\Leftrightarrow\{\dfrac{1}{x_n}\}$是无穷大序列\\
\end{theorem}

\subsubsection{扩充的实数系}
\subsubsection{附录：Stolz定理}
\begin{theorem}
若序列$\{x_n\},\{y_n\}$满足以下条件：
\[(1)0<x_1<x_2<\cdots<x_n<x_{n+1}<\cdots\]
\[(2)\lim_{n\to +\infty}x_n=+\infty\]
如果存在有穷极限【事实上，即使是无穷极限也是可以的，因为只要把分子分母倒过来就变成了有穷极限的0】\\
\[\lim\frac{y_n-y_{n-1}}{x_n-x_{n-1}}=a\]
那么一定也有
\[\lim\frac{y_n}{x_n}=a\]
\end{theorem}
\subsection{函数的极限}
\subsubsection{函数极限的序列式定义}
设$a,A\in\overline{\mathbb{R}},$函数$f(x)$在$a$附近的一个邻域$U_0(a)$有定义\\
如果对一切序列$\{x_n\}\subset U_0(a)$，$x_n\to a$,相应的函数序列值$\{f(x_n)\}$都以A为极限，那么我们就说$x\to a $时，$f(x)$的极限为$A$。

\subsubsection{函数极限的$\varepsilon-\delta$式定义}
设$a,A\in\overline{\mathbb{R}},$函数$f(x)$在$a$附近的一个邻域$U_0(a)$有定义\\
如果对$\forall \varepsilon >0,\exists \delta>0,s.t. \forall x\in U_0(a,\delta)$有
\[|f(x)-A|<\varepsilon\]
以下时一些推广过来的定理具体说明略，只是将
\[\exists N\in \mathbb{N},\forall n>N\]
修改为
\[\exists \delta>0,\forall x\in U_0(a,\delta)\]
\paragraph{夹逼定理\\}
\paragraph{比较原理\\}
若
\[\lim_{x\to a} f(x)<\lim_{x\to a} g(x)\]
则在某个邻域$U_0(a,\delta)$内有
\[f(x)<g(x)\]

\paragraph{推论\\}
则在某个邻域$U_0(a,\delta)$内有
\[f(x)\leq g(x)\]
\[\lim_{x\to a} f(x)\leq \lim_{x\to a} g(x)\]
\paragraph{柯西收敛原理\\}
那么我们就说$x\to a $时，$f(x)$的极限为$A$。
\subsection{单侧极限}
只需要将邻域改为单侧邻域即可\\
\begin{theorem}
这个定理经常被忽视\\
设函数$f(x)$在开区间$(a-\eta,a)$上递增，则
\[\lim_{x\to a-}=\sup_{x\in (a-\eta,a)}f(x)\]
设函数$f(x)$在开区间$(a,a+\eta)$上递增，则
\[\lim_{x\to a+}=\inf_{x\in (a,a+\eta)}f(x)\]
\end{theorem}
\paragraph{求$x$使$F(x)=0$的方法（迭代法）\\}
令$F(x)=f(x)-x$\\
即求$f(x)=x$\\
我们考虑构造数列寻找不动点$f(x^*)=x^*$\\
$x_0$\\
$x_1=f(x_0)$\\
$x_2=f(x_1)$\\
若$x_n$收敛，那么有$x_n\to x^*$
要求：\\
${x_n}$收敛\\
$f(x)$连续\\

\newpage
\section{连续函数}
\subsection{连续与间断}
\paragraph{连续的定义\\}
序列式定义\\
若函数$f(x)$在$x_0$的邻域$U(x_0,\eta)$上有定义，如果对于任何序列$\{x_n\}$满足$x_n\to x_0$，都有\\
\[\lim_{x\to x_0}f(x)=f(x_0)\]
则称函数$f(x)$在$x_0$处连续\\


$\varepsilon-\delta$定义\\
若函数$f(x)$在$x_0$的邻域$U(x_0,\eta)$上有定义，如果\\
\[\forall \varepsilon>0,\exists \delta>0,s.t. \forall x\in U(x_0,\delta),|f(x)-f(x_0)|<\varepsilon\]
则称函数$f(x)$在$x_0$处连续\\

函数$x_0$处连续$\Leftrightarrow$函数在$x_0$处左连续+右连续\\
\paragraph{两类间断点\\}
第一类间断点:
\[\lim_{x\to x_0-}f(x) = \lim_{x\to x_0+} f(x)\neq f(x_0)\]
或\[\lim_{x\to x_0+}f(x)\neq \lim_{x\to x_0-}f(x)\]

\paragraph{小定理\\}
若函数$f(x)$在区间$(a,b)$上单调，则$f(x)$在$(a,b)$上至多只有第一类间断点\\

第二类间断点：
$f(x)$的左右极限至少有一个不存在\\
\begin{theorem}（局部有界性）\\

设函数$f$在$x_0$处连续，则存在$\delta$使得函数在$U(x_0,\delta)$上有界\\
\end{theorem}
\begin{theorem}
若函数$f(x)$在$x_0$出连续，则函数$|f(x)|$在$x_0$处也连续\\
\end{theorem}
\begin{theorem}
设函数$f(x),g(x)$在$x_0$处连续，若$f(x_0)<g(x_0)$则存在$\delta>0$，使得$\forall x\in U(x_0,\delta)$\\
\[f(x)<g(x)\] 
特殊情形体现为【局部保号性】\\
如果$f(x)$是常值函数$f(x)\equiv A,g(x_0)=A$，则存在$\delta>0$，使得$\forall x\in U(x_0,\delta)$\\
\[g(x)>A\] 
\end{theorem}
同样的结果也不赘述
复合函数连续性、函数四则运算的连续性在这里不进行展开\\
\subsection{闭区间上连续函数的重要性质}


\paragraph{引理：闭区间的完备性\\}
闭区间$[a,b]$是完备的，即对于任何序列$\{x_n\}\subset [a,b]$，都有
\[\lim_{n\to +\infty}x_n\in [a,b]\]

\subsubsection{介值定理}
\begin{theorem}(零点存在定理)
若函数$f(x)$在区间$[a,b]$上连续，且$f(a)f(b)<0$，则存在一点$c\in[a,b]$使得
\[f(c)=0\]
\end{theorem}

\begin{theorem}(介值定理)\\
这个定理有两个不同的表述（数院的表述与此不同，并基于\textbf{最大值与最小值定理}）\\
工院教材中的表述：\\
若函数$f(x)$在区间$[a,b]$上连续，则对于$\forall \gamma\in(f(a),f(b)),\exists c\in[a,b]$，使得
\[f(c)=\gamma\]

数院教材中的表述：\\
若函数$f(x)$在区间$[a,b]$上连续，则对于$\forall \gamma\in[\inf_{x\in[a,b]}f(x),\sup_{x\in[a,b]}f(x)],\exists c\in [a,b]$
使得
\[f(c)=\gamma\]
\end{theorem}

\subsubsection{最大值与最小值定理}
\begin{theorem}

设$f(x)$在$[a,b]$上连续，则存在$\xi ,\xi'\in[a,b]$，使得
\[f(\xi)=\inf_{x\in[a,b]}f(x),(\xi')=\sup_{x\in[a,b]}f(x)\]
\end{theorem}
注：由于这里是闭区间，所以上确界与下确界一定是可以取得到的，故可以写成$max$和$min$\\

\begin{theorem}（补充定理： 闭区间上的有界性）\\
若函数$f(x)$在$[a,b]$上连续，则$f(x)$在$[a,b]$上有界\\
\end{theorem}

\subsubsection{一致连续性}
\paragraph{定义\\}
若对于函数$f(x)$,在区间$[a,b]$上有
\[\forall \varepsilon>0,\exists \delta >0,\forall x,y\in[a,b],|x-y|<\delta,|f(x)-f(y)|<\varepsilon\]
那么我们称$f(x)$在$[a,b]$上\textbf{一致连续}\\
\begin{theorem}(一致连续性定理)

若函数$f(x)$在$[a,b]$上连续，则其在$[a,b]$上一致连续\\
\end{theorem}
\subsubsection*{附录：一致连续性的序列式描述}

\begin{theorem}

【这个定理在数院的教材上也有，非常重要！！！】\\
设$E$是$\mathbb{R}$的一个子集，函数$f$在$E$上有定义，则$f$在$E$上一致连续的充分必要条件是：对任何满足条件
\[\lim(x_n-y_n)=0\]
的序列$\{x_n\}\subset E,\{y_n\}\subset E$都有
\[\lim(f(x_n)-f(y_n))=0\]
\end{theorem}

\subsection{单调函数，反函数}

区间的性质：连通性——不连通就不是区间\\
\begin{lemma}
$J\subset\mathbb{R}$是区间$\Leftrightarrow\forall \alpha,\beta\in J,\alpha\leq \beta,\forall\gamma ,\alpha\leq\gamma\leq\beta,\gamma\in J$
\end{lemma}
\begin{theorem}
如果函数$f(x)$在区间$I$上连续，则
\[J=f(I)=\{f(x)|x\in I\}\]
也是一个区间\\
\end{theorem}
\begin{theorem}
如果$f$在$I$上单调，则$f(x)$在$I$连续的充分必要条件是$f(I)$是一个区间\\
\end{theorem}
上面两个定理告诉我们什么？\\
定义域是一个区间+连续$\Rightarrow$值域是一个区间\\
定义域、定义域都是区间+单调$\Rightarrow$函数连续\\

\paragraph{反函数的定义\\}
若$f$在区间$I$连续，则$f(I)$也是一个区间，如果$f$在$I$上还是单调的，则$f$是从$I$到$f(I)$的一个一一对应。\\
同理，对于$\forall y\in f(I)$，存在唯一一个$x\in I$与之对应，因而我们可以在$J=f(I)$定义$g,g(y)$的值为$y$对应的唯一的$x\in I$
这样的函数称为函数$f$的\textbf{反函数}
\[g=f^{-1}\]
\begin{theorem}

若函数$f$在区间$I$上严格单调且连续，则其反函数$g=f^{-1}$在$f(I)$上严格单调且连续\\
\end{theorem}

\paragraph{一些注意点\\}
有反函数未必一定要连续——只需要有一一对应即可\\
同样也未必需要单调\\
\subsection{指数函数与对数函数，初等函数连续性问题小结}

我们省略对指数函数连续性的证明\\
但保留证明中重要的引理\\
\begin{lemma}
对于$\forall a\in \mathbb{R} a>1,\forall p,q\in \mathbb{Q}, |p-q| < 1$有
\[|a^p-a^q|<|p-q|(a-1)a^q\]
\end{lemma}
\begin{theorem}

初等函数在其定义域内都是连续的\\
注：由基本初等函数经过有限次四则运算和复合得到的函数称为初等函数\\
注：基本初等函数包含多项式函数、有理（！）分式函数、三角函数、反三角函数、指数函数、对数函数\\
\end{theorem}
一个重要的极限式子转换
\[u(x)^{v(x)}=e^{v(x)\ln(u(x))}\]
由于指数函数的连续性，我们可以知道
\[\lim u(x)^{v(x)}=e^{\lim(v(x)\ln(u(x)))}\]

\subsection{无穷小量的比较，几个重要的极限}

\paragraph{等价无穷小(无穷大)\\}
若$\lim_{x\to a}\dfrac{\varphi(x)}{\psi(x)}=1$\\

\paragraph{同阶无穷小(无穷大)\\}
若$\lim_{x\to a}\dfrac{\varphi(x)}{\psi(x)}=l$($l$是实数)\\
\paragraph{高阶无穷小(低阶无穷大)\\}

若$\lim_{x\to a}\dfrac{\varphi(x)}{\psi(x)}=0$\\

\paragraph{常见无穷小\\}
以下是在$x\to 0$的情况
\[
\begin{aligned}
&(1) \sin x\thicksim x\\
&(2) \tan x\thicksim x\\
&(3) e^x-1\thicksim x\\
&(4) \ln(x+1)\thicksim x\\
&(5) (1+x)^\mu\thicksim 1+\mu x\\
\end{aligned}
\]
\paragraph{一道习题\\}
请问是否存在$(-\infty,+\infty)$上的连续函数，使得$f(f(x))=e^{-x}$成立？证明你的结论。\\
\begin{lemma}
若定义在$(-\infty,+\infty)$上的函数$f(x)$使得$f(f(x))$有唯一的不动点，则函数$f(x)$也存在唯一的不动点\\
证明：记$f(f(x))$的不动点为$x_0$，即有$f(f(x_0))=x_0$,\\
令$y_0=f(x_0)$，则$f(f(y_0))=f(f(f(x_0)))=f(x_0)=y_0$\\
则$y_0$也为函数$f(f(x))$的不动点\\
由不动点的唯一性知$y_0=x_0$，即有$f(x_0)=x_0$\\
故函数$f(x)$有不动点$x_0$\\
假设函数$f(x)$还有其他不动点$x_1$，则$f(f(x_1))=f(x_1)=x_1$，说明$x_1\neq x_0 $也是函数的不动点，这与题意矛盾\\
故命题成立\\
\end{lemma}
下面我们利用这个引理证明本题\\

记$g(x)=f(f(x))-x=e^{-x}-x$，显然$g(x)$在$\mathbb{R}$上单调递减，且
\[\lim_{x\to -\infty}g(x)=+\infty,\lim_{x\to +\infty}g(x)=-\infty\]
由介值定理知，存在唯一的$x_0$使得$g(x)=0$即$f(f(x))=x$成立\\
同时，根据引理，我们知道
\[\exists ! x_0 \in\mathbb{R},s.t. f(x)=x\]

我们首先考虑$x>x_0$的情况:\\
若$\exists x_1,x_2,s.t.f(x_1)>x_0,f(x_2)<x_0$，则由连续性及介值定理知
\[\exists \eta\in(\min\{x_1,x_2\},\max\{x_1,x_2\}),s.t. f(\eta)=x_0\]\[i.e.\quad e^{-\eta}=f(f(\eta))=f(x_0)=x_0=e^{-x_0}\]

这与指数函数单调性矛盾，故$\forall x>x_0$
\[f(x)>x_0\quad and \quad f(x)<x_0\]
只有一者始终成立。若
\[\forall x>x_0,f(x)>x_0,then,\forall x>x_0,e^{-x}=f(f(x))>x_0=e^{-x_0}\]
这与指数函数单调性矛盾\\
\[\therefore\forall x>x_0,f(x)<x_0\]

那么对于$x<x_0$的情况，与上面的过程相同可以证明$\forall x<x_0$只可能有$f(x)>x_0$\\

我们重新回来考虑$x>x_0$的情况，对于仅剩的可能：对$\forall x>x_0$有$f(x)<x_0$成立\\
我们取$x=x_3>x_0$则有$f(x_3)<x_0$\\
继而根据$x<x_0$时仅剩的可能$f(f(x_3))>x_0$，即
\[e^{-x_3}=f(f(x_3))>x_0=e^{-x_0}\]
这仍然与指数函数的单调性矛盾\\
至此，我们已经证明了$x>x_0$的任何情况都存在矛盾，故证明了这样的函数$f(x)$是不存在的\\
\newpage
\section{导数}
\subsection{导数和微分的概念}
\begin{identification}
(导数的定义)

如果$f$在$x_0$附近有定义,且有穷极限
\[\dfrac{f(x) - f(x_0)}{ x - x_0}\]
存在,则称$f(x)$在$x_0$处可导

事实上,写成
\[f'(x_0) = \lim_{h\to 0} \dfrac{f(x_0 + h) - f(x_0)}{h}\]
更为简单
\end{identification}

\begin{theorem}
  设函数$f,g$在$x_0$处可导,则$f + g$,$cf$在$x_0$也可导
\end{theorem}

\begin{identification}
(单侧导数的定义)

设函数$f$在$\left(x - \eta,x \right]$上有定义,则存在有穷极限
\[\lim_{h\to 0-}\dfrac{f(x + h) - f(x)}{h}\]
则称函数$f$在$x$处左侧可导,同时我们把上述极限成为$f$在$x$处的左导数

同理我们可以定义右侧导数
\end{identification}

\begin{theorem}
  函数$f$在$x$处可导的充分必要条件是其在该点两个单侧导数均存在且相等,即
  \[f'_-(x) = f'_+(x)\]
  此时就有
  \[f'(x) = f'_-(x) = f'_+(x)\]
\end{theorem}

\begin{identification}
(可微的定义)

设函数$f$在$x$附近有定义,若
\[f(x + h) - f(x) = Ah + o(h)\]
其中$A$与$h$无关(但可以与$x$有关,则称$f$在$x$处可微)
\end{identification}

显然我们有$f'(x) = A$

\begin{theorem}
  $f$在$x$可导当且仅当其在$x$处可微
\end{theorem}

\begin{theorem}
  若$f$在$x$处可导,则其在该处连续
\end{theorem}

注意区分$dy$与$\Delta y$!

其中$\Delta y = f(x + h) - f(x)$,成为$f$在$x$处的增量(差分)

而$dy = f'(x) dx = f'(x) \Delta x$,成为$f$在$x$处的微分
\subsection{求导法则,高阶导数}
\subsubsection{基本的求导法则}
\begin{theorem}
  四则运算的求导法则(略)
\end{theorem}

\begin{theorem}
  复合函数的求导法则(略)
\end{theorem}

\subsubsection{一阶微分的形式不变性}
对于$y=f(x)$,我们有$dy = d(f(x)) = f'(x) dx$,\\
假设$x$是因变量且有$x=g(u)$,则我们有
\[dy = d(f(g(u))) = f'(g(u)) g'(u)du = f'(x) dx\]

我们称一阶微分的这种性质为其形式不变性,然而高阶微分并没有这样的形式不变性\\
对于$y=f(x)$,我们有$d^2y = d(f'(x) dx) = f''(x) d^2x$,\\
假设$x$是因变量且有$x=g(u)$,则我们有
\[
\begin{aligned}    
d^2y &= d(f'(g(u)) g'(u)du) = d(f'(g(u)) g'(u))du \\
&= f''(g(u))g'^2(u) du^2 + f'(g(u)) g''(u) du^2 \\
&= f''(x) dx^2 + f'(x) d^2 x\\
\end{aligned}
\]

\paragraph{幂-指数形式的导数}
\[v(x)^{u(x)} = u(x)v(x)^{u(x) - 1} u'(x)+ \ln v(x) \cdot v(x)^{u(x)}  v'(x) \]
相当于分别把这个式子看成指数和幂函数求导加起来

\subsubsection{反函数的求导法则}

\begin{theorem}
  设函数$y = \varphi(x)$在包含$x_0$的区间$I$上\textbf{严格单调且连续},且在$x_0$处可导并且有$\varphi'(x_0) \neq 0$,则$x = \psi(y)$在$y_0 = \varphi(x_0)$处可导,且有
  \[\psi'(y_0) = \dfrac{1}{\varphi'(x_0)} = \dfrac{1}{\varphi'(\psi(y_0))}\]
\end{theorem}

\subsubsection{参数和隐函数方程的求导}
\begin{theorem}
  设有参数表达式$x = \varphi(t),y = \psi(t)$,如果$\varphi,\psi$在$I$上连续,在$t = t_0$处可导,且有$\varphi(t_0) \neq 0$,则有
  \[\dfrac{dy}{dx} = \dfrac{\varphi'(t_0)}{\psi'(t_0)}\]
\end{theorem}

\paragraph{极坐标下极轴与切线的夹角\\}
对于极坐标曲线$r = r(\theta)$,$\text{记}\beta\text{为夹角}$

\[\tan \beta = \dfrac{r(\theta)}{r'(\theta)}\]

\paragraph{隐函数求导\\}

\begin{example}
  对于$x^2 + y^2 = 1$,以$x$为自变量,把$y$看作关于$x$的函数,两边对$x$求导得到
  \[2x + 2y y' = 0\]
\end{example}

\subsubsection{高阶导数}

\begin{theorem}
(莱布尼茨公式)  

设$u,v$在$x_0$处$n$阶可导,则乘积$uv$在$x_0$处同样$n$阶可导,则有
\[(uv)^{(n)} = \sum_{k = 0}^n \binom{n}{k} u^{(n - k) v^{(k)}}\]
\end{theorem}

\paragraph{易错点}
对于反函数求导
\[G'(y) = \dfrac{1}{F'(G(y))}\]
\[
\begin{aligned}
  G''(y) &= -\dfrac{1}{\parameter{F'(G(y))}^2} F''(G(y)) G'(y)\\
  &=-\dfrac{F''(G(y))}{\parameter{F'(G(y))}^2}  \dfrac{1}{F'(G(y))}\\
  &=-\dfrac{F''(G(y))}{\parameter{F'(G(y))}^3}\\
\end{aligned}  
\]

同样,对于参数方程(符号沿用之前的)
\[\dfrac{dy}{dx} = \dfrac{\psi'(t)}{\varphi'(t)}\]
\[
\begin{aligned}
  \dfrac{d^2y}{dx^2} &= \dfrac{d\parameter{\frac{dy}{dx}}}{dx}\\
  &=\dfrac{\parameter{\frac{\psi'(t)}{\varphi'(t)}}'}{\varphi'(t)}\\
  &=\dfrac{\psi''(t)\varphi'(t) -\psi'(t) \varphi''(t)}{\parameter{\varphi'(t)}^3}
\end{aligned}  
 \]

更要注意的是$\dfrac{d^2 y}{dx^2} \neq \dfrac{dy}{dx} \cdot \dfrac{dy}{dx},\text{并且}\dfrac{d^2 y}{dx^2} \neq \dfrac{d^2 y}{dt^2} \cdot \dfrac{dt}{dx}\cdot \dfrac{dt}{dx}$

\subsection{无穷小增量公式和有限增量公式}
无穷小增量公式
\[f(x + h) = f(x) + f'(x) h + o(h)\quad\quad(h\to 0)\]

\paragraph{应用\\}
主要是近似计算

\begin{example}
  求$\sqrt[5]{270}$的近似值
  \[f(x + \Delta x) = f(x) + f'(x) \Delta x + o(\Delta x)\]
\end{example}
有限增量公式(这个依赖于后面我们讲的中值定理)

\[f(x) = f(x_0) + f'(\xi) (x_0 - x)\]
或者写作
\[f(x + h) = f(x) + f'(x +\theta h) h\]

\paragraph{应用\\}
主要是一些证明
\subsection{重要定理补充}
\begin{theorem}（重要补充定理：Darboux定理）\\ 
若 $f(x)$在 $[a,b]$上可导，且有 $f'(a)f'(b) < 0 $ 则 $\exists \xi \in (a,b)$使得函数
\[f'(\xi) = 0\]
\end{theorem}
证明：不妨假设$f'(a) > 0, f'(b) < 0$ 则存在 $x_1 > a , x_2 < b , s.t. f(x_1) > f(a) , f(x_2) > f(b)$

在闭区间$[a,,b]$上，利用最大值与最小值定理，存在一个$\xi\in [a,b]$使得 
\[f(\xi) = \max_{x\in [a,b]}f(x)\]

又由于$f(\xi) > f(x_1) > f(a), f(\xi) > f(x_2) > f(b)$ ，从而我们知道$\exists \xi\in (a,b),s.t. f(\xi)$是区间$[a,b]$上的最大值，从而便是极大值，便有$f'(\xi) = 0$\\

拓展形式：
若 $f(x)$在 $[a,b]$上可导，且 $c$为介于$f'(a)$ 与 $f'(b)$之间的值，那么存在 $\xi \in (a,b),s.t.$
\[f'(\xi) = c\]

证明：\\
①若$f'(a) = f'(b)$ 显然成立（是否要把条件中的开区间改成闭区间？）\\
②若$f'(a) > f'(b)$ 那么对于 $\forall \eta\in (f'(b),f'(a))$，我们构造一个 $F(x) =  f(x) - \eta x$
从而
\[F'(a) = f'(a) - \eta < 0 , F'(b) = f'(b) - \eta > 0 \]
化归为我们已经证明的情况。\\
③若$f'(a) < f'(b)$，情况同②，在这里我们也省略\\


这个定理看起来和介值定理及其推广没有什么区别，但是事实上有很大的差异：介值定理是基于函数本身连续的基础上的，然而Darboux定理并没有要求导函数的任何相关性质！\\

\paragraph{应用\\} 
作为Darboux重要性的事例，下面我们来证明导函数的一个重要性质：\\

导数只可能有第二类间断点，不可能有第一类间断点。

证明采用反证，需要使用Darboux定理\\
首先我们可以说明第一类间断点是不可能存在的，因为这必然会导致在间断点的附近，不满足Darboux定理的相关条件（即介于两个值导数值之间的其他导数值无法被取到，这是矛盾的）\\
其中第二类间断点也可以分为两类，第一类是震荡型，另一类是单调型，\\
震荡型的函数可以为
\begin{equation}
f(x) =  
\begin{cases}
x^{2}\sin \dfrac{1}{x} &,x\neq 0\\
0 &, x=0\\
\end{cases}
\end{equation}
这个函数在$x=0$处可导，因为
\[\lim_{h\to 0}\dfrac{f(h+0)-f(0)}{h}=\lim_{h\to 0}h\sin\frac{1}{h} = 0\]
显然这个函数在其他地方可导，从而，$f(x)$可导。\\

而这个函数的导数可以表示为
\[f'(x) = 2x\sin \frac{1}{x} - \cos\frac{1}{x}\]
显然$x=0$处导数不连续，并且应当为震荡间断点\\

单调型是不可能存在的，同样可以采用Darboux定理\\

一个容易混淆的例子是
\[f(x)=\sqrt[3]{x}\]
事实上这个函数在$x=0$处是不可导的，因为
\[\lim_{h\to 0 }\dfrac{f(h+0)-f(0)}{h}=\lim_{h\to 0} \dfrac{1}{\sqrt[3]{h^2}}\]
这个极限显然是不存在的\\

而我们刚才讨论的前提是“导函数存在”\\

事实上，关于这个问题的讨论告诉我们，一个函数$f(x)$其导函数存在（或者函数可导），与导函数的连续性没有直接的联系，导函数存在，不意味着导函数的连续（尽管我们能肯定其不存在第一类间断点和第二类间断点中的单调（到无穷的）情况）\\

\begin{theorem} （重要定理补充：导数极限定理）\\

设函数$f$在点 $x_0$的某邻域 $U(x_0)$内连续 ，在 $U_0(x_0)$内可导，且极限$\lim_{x\to x_0}f'(x)$存在，则$f$在$x=x_0$处可导，且$f'(x_0) = \lim_{x\to x_0}f'(x)$\\
\end{theorem}

证明：\\
分别按照左右导数给出证明：\\
我们先考虑右导数：\\
任取$x\in U_0^{+}(x_0)$，$f$在$[x_0,x]$上满足Lagrange定理的条件，从而有
\[\exists \eta \in (x_0,x),s.t.\dfrac{f(x)-f(x_0)}{x-x_0}=f'(\eta)\]
由于$x_0 < \eta < x$，所以当$x\to x_0^+ $时有$\eta \to x_0^+$，从而我们对上式取极限
\[\lim_{x\to x_0^+}\dfrac{f(x)-f(x_0)}{x-x_0} = \lim_{x\to x_0^+} f'(\eta)= f'(x_0+0)\]
弄清楚这个式子的意思！！\\
等式左边是$x=x_0$这点的右导数，等式右边是导函数在$x_0$处的右极限\\
同理可以证得左极限的情况，因而我们可以证明我们想要的结论\\

\section{一元微积分定理的再补充}
\subsection{微分中值定理}
\begin{theorem}（罗尔（Rolle）定理）

  设$f(x)$在$[a,b]$上连续,在$(a,b)$上可导,且满足
  \[f(a) = f(b)\]
  则存在$\xi\in(a,b)$,使得
  \[f'(\xi) = 0\]
\end{theorem}

\begin{theorem}（拉格朗日（Langrange）中值定理）
    
  设$f(x)$在$[a,b]$上连续,在$(a,b)$上可导,则存在$\xi\in(a,b)$,使得
  \[f'(\xi) = \dfrac{f(b)-f(a)}{b-a}\]

\end{theorem}

\begin{theorem}（柯西中值定理）
        
  设$f(x),g(x)$在$[a,b]$上连续,在$(a,b)$上可导,则存在$\xi\in(a,b)$,使得
  \[\dfrac{f'(x)}{g'(x)} = \dfrac{f(b)-f(a)}{g(b)-g(a)}\]

\end{theorem}

\subsection{洛必达法则}
更多用于理论的分析而不是计算\\
\begin{theorem}
若函数在$f,g$在邻域$U_0(a,\delta)$上可导，若满足下列条件
\[
\begin{aligned}
&(1)\lim_{x\to a}f(x) =\lim_{x\to a } g(x) = 0\\
&(2)g(x) \neq 0 ,\forall x\in U_0(a,\delta)\\
&(3)\dfrac{f'(x)}{g'(x)} = l(l\text{为有限数},\pm\infty,\infty)\\
\end{aligned}
\]
则
\[\lim_{x\to a}\dfrac{f(x)}{g(x)}=\lim_{x\to a}\dfrac{f'(x)}{g'(x)}\]
\end{theorem}

\begin{theorem}
若函数在$f,g$在邻域$U = \{|x| > a > 0\}$上可导，若满足下列条件
\[
\begin{aligned}
&(1)\lim_{x\to \infty}f(x) =\lim_{x\to \infty } g(x) = 0\\
&(2)g(x) \neq 0 ,\forall x\in U\\
&(3)\dfrac{f'(x)}{g'(x)} = l(l\text{为有限数},\pm\infty,\infty)\\
\end{aligned}
\]
则
\[\lim_{x\to \infty}\dfrac{f(x)}{g(x)}=\lim_{x\to \infty}\dfrac{f'(x)}{g'(x)}\]
\end{theorem}

\begin{theorem}\label{lh3}
若函数在$f,g$在邻域$U_0(a,\delta)$上可导，若满足下列条件
\[
\begin{aligned}
&(1)\lim_{x\to a } g(x) = \infty\\
&(2)g(x) \neq 0 ,\forall x\in U_0(a,\delta)\\
&(3)\dfrac{f'(x)}{g'(x)} = l(l\text{为有限数},\pm\infty,\infty)\\
\end{aligned}
\]
则
\[\lim_{x\to a}\dfrac{f(x)}{g(x)}=\lim_{x\to a}\dfrac{f'(x)}{g'(x)}\]
\end{theorem}

\begin{theorem}\label{lh4}
若函数在$f,g$在邻域$U = \{|x| > a > 0\}$上可导，若满足下列条件
\[
\begin{aligned}
&(1)\lim_{x\to \infty } g(x) = \infty\\
&(2)g(x) \neq 0 ,\forall x\in U\\
&(3)\dfrac{f'(x)}{g'(x)} = l(l\text{为有限数},\pm\infty,\infty)\\
\end{aligned}
\]
则
\[\lim_{x\to \infty}\dfrac{f(x)}{g(x)}=\lim_{x\to \infty}\dfrac{f'(x)}{g'(x)}\]
\end{theorem}
我们要注意一点,~\ref{lh3}和~\ref{lh4}并不依赖分子是无穷,就算其是非无穷极限也可以\\
\subsection{泰勒公式}
若函数$f(x)$在$x_0$处$n$阶可导，则
\[f(x) = f(x_0) + f'(x_0) (x-x_0) + \dfrac{f''(x_0)}{2 !} (x-x_0)^2 + \cdots + \dfrac{f^{(n)}(x_0)}{n !} (x-x_0)^n +o ((x-x_0)^n)\]
\paragraph{皮亚诺余项\\}
形如$o ((x-x_0)^n)$的项称为Peano余项,可以帮助我们估计某点处的极限值

我们可以利用这一点解决很多极限类的问题,我们可以记忆以下常见的展开(在$x = 0$处展开)(主要是几个不太常见的)
\[\tan  x = x + \dfrac{x^3}{3} + \dfrac{2x^5}{15} + o(x^5) \]

\[\arcsin x = x + \dfrac{x^3}{6} + \dfrac{3x^5}{40} + o(x^5)\]

\[\arccos x = \dfrac{\pi}{2} - \arcsin x\]

\[\arctan x = x - \dfrac{x^3}{3} + \dfrac{x^5}{5} + o(x^5)\]

\paragraph{拉格朗日余项}
若函数$f(x)$在$[a,b]$处$n$阶可导，在$(a,b)$上$(n+1)$阶可导，则
\[f(x) = f(x_0) + f'(x_0) (x-x_0) + \dfrac{f''(x_0)}{2 !} (x-x_0)^2 + \cdots + \dfrac{f^{(n)}(x_0)}{n !} (x-x_0)^n + \dfrac{f^{(n+1)}(\xi)}{(n+1) !} (x-x_0)^{n+1}\]
\[\forall x,x_0\in (a,b),\xi\text{介于}x_0,x\text{之间}\]

证明思路： 构造余项函数证明这个余项函数是一个无穷小量，由于泰勒公式的唯一性即可，利用柯西中值定理\\

提示:构造\[F(t) = f(x) - (f(t) + f'(t) (x-t) + \dfrac{f''(t)}{2 !} (x-t)^2 + \cdots + \dfrac{f^{(n)}(t)}{n !} (x-t)^n )\]\[G(t) = (x-t) ^ {n+1}\]
同样对于其中的$\xi$可以表示为$\xi = x_0 + (x- x_0) \theta$\\

推论：如果$f^{(n+1)}(x)\equiv 0,\forall x\in (a,b)\Rightarrow$ $f$为$n$次多项式\\


当$x=0$时称为带Lagrange余项的麦克劳林公式\\

一般记忆方式（就是在原先$n$次的基础上再写一项，只是把$x_0$换为$\xi$）\\
核心是要记牢\textbf{高阶导数}
\[
\begin{aligned}
    (e^x)^{(n + 1)} &= e^x\\
    &\\
    (\ln(1 + x))^{(n + 1)} &= (-1)^n \dfrac{n !}{(1 + x)^{n + 1}} \\
    &\\
    (\sin x)^{(2n + 1)} &= \sin \parameter{x + (2n + 1) \frac{\pi}{2}} = (-1)^{x} \cos x\\
    &\\
    (\cos x)^{(2n + 2)} &= \cos \parameter{x + (2n + 2) \frac{\pi}{2}} = (-1)^{n + 1} \cos x\\   
    &\\
    ((1 + x)^\alpha)^{(n + 1)} &= \alpha(\alpha - 1)\cdots (\alpha - n)(1 + x)^{\alpha - n - 1}\\
\end{aligned}   
\]

\paragraph{一些需要记忆的拉格朗日余项展开形式}
\[
\begin{aligned}    
e^x &= 1 + x + \dfrac{x^2}{2!} + \cdots + \dfrac{x^n}{n!} + \dfrac{e^{\theta x} x^{n+1}}{(n+1)!}\\
&\quad\quad\quad\quad(-\infty < x + \infty , 0 < \theta < 1)\\
\ln(1 + x) &= x - \dfrac{x^2}{2} + \cdots + (-1)^{n-1} \dfrac{x^n}{n} + (-1)^n \dfrac{x^{n+1}}{(n + 1)(1 + \theta x)^{n + 1}}\\
&\quad\quad\quad\quad(-1 < x < 1 , 0 < \theta < 1)\\
\sin x &= x - \dfrac{x^3}{3!} +\cdots + (-1)^{n-1} \dfrac{x^{2n-1}}{(2n-1)!}+(-1)^{n}\dfrac{\cos\theta x}{(2n + 1)!}x^{2n+1}\\
&\quad\quad\quad\quad(-\infty < x + \infty , 0 < \theta < 1)\\
\cos x &= 1 - \dfrac{x^2}{2!} + \cdots + (-1)^n \dfrac{x^2n}{(2n)!} + (-1)^{n+1} \dfrac{\cos \theta x}{(2n + 2)!} x^{2n + 2}\\
&\quad\quad\quad\quad(-\infty < x + \infty , 0 < \theta < 1)\\
(1 + x)^\alpha &= 1 + \alpha x + \dfrac{\alpha(\alpha - 1)}{2!} x^2 + \cdots + \dfrac{\alpha \cdots (\alpha - n + 1)}{n !} x^n +\dfrac{\alpha \cdots (\alpha - n)(1 + \theta x)^{\alpha - n -1}}{(n + 1)!} x^{n + 1}\\
&\quad\quad\quad\quad(-\infty < x + \infty , 0 < \theta < 1)\\
\end{aligned}
\]

\paragraph{泰勒级数}

\[
f(x) = \lim_{n\to +\infty} \sum_{k = 0} ^n \dfrac{f^{(k)}(a)}{k!}(x-a)^k
= \sum_{k}^{+\infty} \dfrac{f^{(k)}(a)}{k!}(x-a)^k
\]
称为Taylor级数，$R_{2n+1}\to 0$取决于$a$，$x$的范围称为收敛半径\\

注意$\ln(1+x)$的收敛半径不是$\mathbb{R}$，而是$0 \leq x <1$ ，这点与其他大多数基本初等函数不同\\
如果在收敛半径内我们可以写成泰勒级数，否则不可以\\

若$n\to +\infty,R(n) \to 0$，则$f(x)$展开成Taylor级数，称为是解析的\\

无穷次可导但不解析
\begin{equation}
    f(x) = 
    \begin{cases}
        e^{-\dfrac{1}{x}}, & x > 0\\
        0, & x \leq 0
    \end{cases}
\end{equation}
    
$f$无穷次可导 $f^{(k)}(0) = 0,\forall k\in \mathbb{N}$\\
\[\Rightarrow \sum_{k=0}^n\dfrac{f^{(k)}(0)}{k!}x^k = 0\]
\[f(x)\neq \sum_{k=0}^{+\infty}\dfrac{f^{(k)}(0)}{k!}x^k\]

\subsection{用导数研究函数}
\subsubsection{极值}
\begin{identification}
  (内点)

设$I$是一个区间$x_0\in I$,如果存在$\eta >0$,使得$U(x_0,\eta) \subset I$,则称$x_0$为$I$的一个内点

区间$I$去除(可能存在的)端点后剩下的点都是内点,它们组成的集合是一个开区间,记为$I^0$\\
\end{identification}

\begin{identification}
设函数$f$在$I$上有定义,$x_0\in I^0$,若存在一个邻域$U(x_0,\delta)\subset I$使得$\forall x \in U(x_0,\delta)$都有
\[f(x_0) \geq f(x) \quad \parameter{f(x_0) \leq f(x)}\]
则称$f$在$x_0$处取极大值(极小值)

若存在一个邻域$U_0(x_0,\delta)\subset I$使得$\forall x \in U_0(x_0,\delta)$都有
\[f(x_0) > f(x) \quad \parameter{f(x_0) < f(x)}\]
则称$f$在$x_0$处取严格极大值(严格极小值)
\end{identification}

\begin{theorem}
  (极值的必要条件)

  设函数$f$在$I$上有有定义,$x_0\in I$,且$f$在$x_0$处取到极值,若$f$在$x_0$处可导,则有
  \[f'(x_0) = 0\]
\end{theorem}

注意:如果$f$在$x_0$是不可导的话,如$y = |x|$,那么上述条件也不是必要的

\begin{identification}
我们把使$f'(x) = 0$的点称为$f$的临界点(驻点)
\end{identification}
\begin{theorem}
  设$f(x)$在$[a,b]$上连续,在$(a,b)$上可导,若$f'(x) = 0$在$(a,b)$上只有有限个根$x_1,\cdots ,x_n$,则函数$f$在$[a,b]$上的最值可以表示为
  \[M = \max\{f(a),f(x_1),\cdots,f(x_n),f(b)\}\]

  和
  \[m = \min\{f(a),f(x_1),\cdots,f(x_n),f(b)\}\]
\end{theorem}

\begin{theorem}（极值的第一充分条件）\\
$f(x)$在$U(x_0,\delta)$上可导且导函数连续，有以下条件\\
(1)若$\forall x\in a-\delta<x<a,f'(x) > 0$,$\forall x\in a<x<a + \delta,f'(x) < 0$则$f(x)$在$x_0$处取极大值\\
(2)若$\forall x\in a-\delta<x<a,f'(x) < 0$,$\forall x\in a<x<a + \delta,f'(x) > 0$则$f(x)$在$x_0$处取极小值\\
\end{theorem}

\begin{theorem}（极值的第二充分条件）\\
若函数$f(x)$在区间$U(x_0,\delta)$上二阶可导，则满足以下条件\\
$f'(x_0) = 0$对应临界点（未必是极值点）\\
若$f''(x_0) > 0$ 则$(x)$在$x_0$处取极小值\\
$f''(x_0) < 0$ 对应极小值\\
\end{theorem}

\begin{theorem}（极值的第三充分条件）\\

$f(x)$在$U(x_0,\delta)$上$n$阶可导
\[f'(x_0) = f''(x_0) = \cdots =f^{(n-1)}(x_0) = 0,f^{(n)}(x_0)\neq 0\]
则

(1)$n$为偶数

若$f^{(n)}(x_0) > 0$则$f(x)$在$x_0$处取极小值

若$f^{(n)}(x_0) < 0$则$f(x)$在$x_0$处取极大值

(2)$n$为奇数，$x_0$不是极值点
\end{theorem}

\subsubsection{函数的凹凸性}

\begin{identification}
\[f(tx_1 + (1-t)x_2)  \leq tf(x_1) + (1-t) f(x_2),\forall x_1 ,x_2 \in [a,b]\]
则称函数$f(x)$为$[a,b]$上的凸函数(如果将$\leq$换为$<$则有是严格凸函数)\\

\[f(tx_1 + (1-t)x_2)  \geq tf(x_1) + (1-t) f(x_2),\forall x_1 ,x_2 \in [a,b]\]
则称函数$f(x)$为$[a,b]$上的凹函数(如果将$\geq$换为$>$则有是严格凹函数)\\
\end{identification}

\begin{theorem}
$\forall x_1,x_2 ,x_3 \in [a,b] ,x_1< x_2 < x_3$
\[\Delta = \det
\begin{bmatrix}
    1&x_1&f(x_1)\\
    1&x_2&f(x_2)\\
    1&x_3&f(x_3)\\    
\end{bmatrix}
\begin{aligned}
&\geq 0 \Leftrightarrow f\text{是凸函数}\\
& > 0\Leftrightarrow f\text{是严格凸函数}\\
\end{aligned}
\]
\end{theorem}
这个定理刻画了什么?我们在线性代数中已经学到了这个行列式的值对应着三角形有向面积的两倍,而这个值的正负实际上表征了这三个点是否满足某种逆时针关系,由于凸函数中$(x_2 - x_1 ,f(x_2)-f(x_1))$与$(x_3 - x_1 ,f(x_3)-f(x_1))$恰好满足了这种逆时针关系\\

这个定理广泛地用于各类证明中，如果题目要求我们证明一些相当显然（或者和常用定理类似但又不完全相同的定理（比如不证明导数的单调性而是证明单侧导数的单调性，我们无法直接套用现有定理，但是可以模仿这些定理的推导，利用这个行列式形式的定理））\\
\begin{center}
    \begin{tikzpicture}
      \begin{axis}[
        xlabel=$x$,
        ylabel=$f(x)$,
        axis lines=middle,
        xmin=-0.3, xmax=2,
        ymin=-1, ymax=4,
        width=\textwidth,
        height=8cm,
        xtick = \empty,
        ytick = \empty
      ]
      \addplot[black,line width=1pt,domain=0.5:1.3,samples=200] {tan(deg(x))};
      \addplot[only marks, mark=*] coordinates {(0.6, {tan(deg(0.6))}) (1.2, {tan(deg(1.2))}) (0.9, {tan(deg(0.9))})};

      \addplot[red, line width=1pt] coordinates {(0.6, {tan(deg(0.6))}) (1.2, {tan(deg(1.2))})};

      \addplot[red, line width=1pt] coordinates {(0.6, {tan(deg(0.6))}) (0.9, {tan(deg(0.9))})};
      
      \addplot[red, line width=1pt] coordinates {(0.9, {tan(deg(0.9))}) (1.2, {tan(deg(1.2))})};

      \node at (axis cs:0.5, {tan(deg(0.5))}) [anchor=north west] {$(x_1,f(x_1))$};

      \node at (axis cs:1.2, {tan(deg(1.2) - 0.8)}) [anchor=south west] {$(x_3,f(x_3))$};

      \node at (axis cs:0.85, {tan(deg(0.85) - 0.2)}) [anchor=north west] {$(x_2,f(x_2))$};

      \end{axis}
    \end{tikzpicture}
  \end{center}
\begin{theorem}
\[
\begin{aligned}    
&\forall x_1,x_2,x_3 \in [a,b],x_1<x_2 < x_3,\\
&\dfrac{f(x_2)-f(x_1)}{x_2-x_1}\leq\dfrac{f(x_3)-f(x_2)}{x_3-x_2} \Leftrightarrow f\text{是凸函数}\\
\end{aligned}
\]
\end{theorem}
根据上面一个定理我们很容易给出证明,这里仅仅做一些证明的提示,行列式第三行减去第二行,第二行减去第三行,再按第一列展开,这一些过程都是可逆的,因此都是当且仅当,于是证明了充分必要性\\

\begin{theorem}(凸函数的可导性)
若$f(x)$是$[a,b]$上的凸函数,则$\forall x\in (a,b),f(x)$在$x$处左右可导(没有说该点可导,只是左右可导)且连续\\
\end{theorem}
\begin{cproof}
我们可以证明割线斜率的单调性
\[\dfrac{f(x+h) - f(x)}{h} < \dfrac{f(x_1) - f(x)}{x_1 - x}\]
我们可以把左侧的函数视为一个关于$h$的函数,这个函数单调递增(1)而有上界,因此其极限是存在的,故$f'_+(x)$存在\\

(1)中的单调性我们同样可以利用前面的行列式得到(准确来说，是用第一行和第二行同时减去第三行，我们即可证明这样的单调性)\\

同理左侧的函数的右极限也存在,故$f'_+(x)$存在\\

\begin{center}
    \begin{tikzpicture}
      \begin{axis}[
        xlabel=$x$,
        ylabel=$f(x)$,
        axis lines=middle,
        xmin=-0.3, xmax=2,
        ymin=-1, ymax=4,
        width=\textwidth,
        height=8cm,
        xtick = \empty,
        ytick = \empty
      ]
      \addplot[black,line width=1pt,domain=0.5:1.3,samples=200] {tan(deg(x))};
      \addplot[only marks, mark=*] coordinates {(0.6, {tan(deg(0.6))}) (1.2, {tan(deg(1.2))}) (0.9, {tan(deg(0.9))})};

      \addplot[red, line width=1pt] coordinates {(0.6, {tan(deg(0.6))}) (1.2, {tan(deg(1.2))})};

      \addplot[red, line width=1pt] coordinates {(0.6, {tan(deg(0.6))}) (0.9, {tan(deg(0.9))})};
      
      \addplot[red, line width=1pt] coordinates {(0.9, {tan(deg(0.9))}) (1.2, {tan(deg(1.2))})};

      \node at (axis cs:0.5, {tan(deg(0.5))}) [anchor=north west] {$(x-h,f(x-h))$};

      \node at (axis cs:1.2, {tan(deg(1.2) - 0.8)}) [anchor=south west] {$(x_1,f(x_1))$};

      \node at (axis cs:0.85, {tan(deg(0.85) - 0.2)}) [anchor=north west] {$(x,f(x))$};

      \end{axis}
    \end{tikzpicture}
  \end{center}
这两个条件可以推出$f(x)$在$(a,b)$连续(Why?)\\
我们回顾以下导数的定义中的内容,左导数存在$\Rightarrow$函数左连续,右导数存在$\Rightarrow$函数右连续\\
\end{cproof}
\begin{theorem}
如果函数$f(x)\in C[a,b]$,且$f(x)$在$(a,b)$上可导,则$f(x)$是凸函数当且仅当$f'(x)$在$(a,b)$上单调递增\\

如果函数$f(x)\in C[a,b]$,且$f(x)$在$(a,b)$上可导,则$f(x)$是严格凸函数当且仅当$f'(x)$在$(a,b)$上严格单调递增\\
\end{theorem}
\begin{cproof}
先证明必要性$\Rightarrow$(非严格)
$x_1,x_2 \in(a,b),x_1 < x_2 $,要证明 $f'(x_1) < f'(x_2)$\\
取$h$充分小$h>0$,$x_1-h , x_2 + h \in (a,b)$\\
\[\dfrac{f(x_1) - f(x_1 - h)}{h} < \dfrac{f(x_2) - f(x_1)}{x_2 - x_1} < \dfrac{f(x_2 + h) - f(x_2)}{h}\]
令$h\to 0$\\
\[f'(x_1) \leq \dfrac{f(x_2) - f(x_1)}{x_2 - x_1} \leq f'(x_2)\]
对于严格的证明很简单,只要在$x_1和x_2$之间插入一个$x'$即可,因为两个割线斜率之间保持严格的不等关系,因此即可证明严格的递增\\

下面我们证明充分性$\Leftarrow$\\

由Lagrange中值定理
\[\dfrac{f(x_2) - f(x_1)}{x_2 - x_1} = f'(\xi_1) ,\exists\xi_1 \in (x_1,x_2)\]
\[\dfrac{f(x_3) - f(x_2)}{x_3 - x_2} = f'(\xi_2) ,\exists\xi_2 \in (x_2,x_3)\]

由于$f'(x)$(严格)递增$\Rightarrow f'(\xi_1) (>) \geq f'(\xi_2)\Rightarrow f$是(严格)凸函数\\

\end{cproof}

\begin{theorem}
$f\in C[a,b]在(a,b)$上二阶可导,则$f$是凸函数$\Leftrightarrow f''(x) \geq 0$\\
$f$是严格凸函数 $\Leftrightarrow f''(x) \geq 0$,且$f''(x)$不在$(a,b)$上任意开子区间里恒为零($f(x)$在$(a,b)$上不包含直线段)
\end{theorem}
\begin{cproof}
$f$为凸函数$\Leftrightarrow f'(x)$单调递增 $\Leftrightarrow f''(x) \geq 0$\\

$f$为严格凸函数 $\Leftrightarrow f'(x)$严格单调递增 $\Leftrightarrow f''(x) \geq 0$ 且其不在$(a,b)$中任何开子区间中恒为零\\
\end{cproof}
在某个区间上$f''(x)\geq 0$,则极小值是最小值\\

\begin{theorem}
$f\in C[a,b]$在$(a,b)$上可导,则\\
\[f\text{是(严格)凸函数}\Leftrightarrow \forall x_0 \in (a,b) f(x) (>)\geq f(x_0) + f'(x_0) (x-x_0)\]

这个定理在直观理解上不难,其实就讲了函数与其切线的位置关系,这里只给出证明的一些要点,证明必要性$(\Rightarrow)$时,我们分为$x > x_0$和$x < x_0$两种情况讨论即可,注意正负不同的时候不等号的变号情况\\

证明充分性($\Leftarrow$)时只需证明$f'(x)$(严格)递增\\

只要将结论的条件对于$\forall x_1,x_2\in(a,b)(x_1<x_2)$分别写出来,然后进行等式的变形,即可得到 $f'(x_1) (<)\leq f'(x_2)$\\
\end{theorem}

\subsubsection{函数凹凸性的应用}

\begin{theorem}(Jensen不等式)
$f$在$[a,b]$上是凸函数,$x_i \in [a,b]$, $t_i$(权重) $> 0(i = 1,2,\cdots n) ,\sum_{i = 1} ^{n} t_i = 1$,则\\
\[f(t_1 x_1 + t_2 x_2 + \cdots + t_n x_n) \leq t_1f(x_1) + t_2 f(x_2) + \cdots + t_n f(x_n)\]

当$f$为严格凸函数且$x_i$不全等,则不等号严格成立\\
\end{theorem}
\begin{cproof}
证明利用数学归纳法

$n = 2$时即是凸函数的定义

假设$n = k$时不等式成立

则$n = k+1$时,设$t_i > 0 ,\sum_{i = 1} ^ {k + 1} t_i= 1$,同时设$\lambda_i = \dfrac{t_i}{1- t_{k+1}}$,那么有
\[\lambda_i > 0, \sum_{i = 1} ^k \lambda_i = 1\]
\[
\begin{aligned}
f(t_1 x_1 + t_2 x_2 + \cdots + t_{k+1} x_{k+1}) &=f\parameter{(1-t_{k+1})\dfrac{t_1 x_1 + t_2 x_2 + \cdots + t_k x_k}{1-t_{k+1}} + t_{k+1} x_{k+1}}\\
&\leq (1-t_{k+1})f\parameter{\dfrac{t_1 x_1 + t_2 x_2 + \cdots + t_k x_k}{1-t_{k+1}}} + t_{k+1}f(x_{k+1})\\
&=(1-t_{k+1})f(\lambda_1 x_1 + \cdots + \lambda_k x_k) + t_{k+1}f(x_{k+1})\\
&\leq(1-t_{k+1})(\lambda_1 f(x_1) +\cdots \lambda_k f(x_k)) + t_{k+1}f(x_{k+1})\\
&=(1-t_{k+1})\parameter{\dfrac{t_1}{1-t_{k+1}}f(x_1) +\cdots \dfrac{t_k}{1-t_{k+1}}f(x_k)}+ t_{k+1}f(x_{k+1})\\
&=t_1 f(x_1) + \cdots + t_k f(x_k) + t_{k+1} f(x_{k+1})\\
\end{aligned}
\]
\end{cproof}

这些带有$n$的不等式很多都可以化为Jensen不等式来完成(比如均值不等式链)
基本不等式的证明
\begin{example}
\[f(x)= \ln x,f'(x) = \frac{1}{x} , f''(x) = \frac{-1}{x^2} < 0\]
\[f(x)\text{是凹函数}\]
\[\text{取}x_i>0, t_i = \frac{1}{n},\therefore \ln(\frac{x_1 + x_2 +\cdots + x_n}{n})  \geq\sqrt[n]{x_1 x_2 \cdots x_n}\]

\[f(x)=-\ln x,f'(x) = -\frac{1}{x} , f''(x) = \frac{1}{x^2} < 0\]
\[f(x)\text{是凸函数}\]
\[\text{取}x_i>0, t_i = \frac{1}{n},\therefore -\ln(\frac{\frac{1}{x_1} + \frac{1}{x_2} +\cdots + \frac{1}{x_n}}{n})  \leq -\frac{1}{n}\ln\frac{1}{x_1}-\frac{1}{n}\ln\frac{1}{x_2} - \cdots -\frac{1}{n}\ln\frac{1}{x_n}  \]
\[\frac{\frac{1}{x_1} + \frac{1}{x_2} +\cdots + \frac{1}{x_n}}{n} \leq \sqrt[n]{x_1 x_2 \cdots x_n}\]

\end{example}
此类证明的核心是构造一个合适的具有稳定凹凸性的$f(x)$\\

\begin{example}
  $a_i$不全等$a_i > 0$,证明:\\
  \[f(x) = (\dfrac{a_1^2 + a_2^2 +\cdots + a_n^2}{x})^{\frac{1}{x}}\text{在}(-\infty,+\infty)\text{上递增}\]
首先我木需要补充定义$x=0$时 $f(0) = \lim_{x\to 0}f(x) = \sqrt[n]{a_1\cdots a_n}$\\

这个函数的求导通常应当使用对数求导法\\

\[
\begin{aligned}
(\ln f(x))' &= \dfrac{f'(x)}{f(x)}\\
&= -\dfrac{1}{x^2}[\ln(a_1^2 + a_2^2 + \cdots + a_n^2)/n + \dfrac{a_1^x \ln  a_1+ \cdots + a_n^x \ln a_n}{a_1^x + \cdots + a_n^x}] > 0\\
\end{aligned}
\]
\end{example}
Hoder不等式\\
\begin{example}
  \[a_i > 0, b_i > 0, i = 1, 2, \cdots, n\]
  \[\sum_{i = 1}^n a_i b_i \leq (\sum_{i= 1}^n a_i^p)^\frac{1}{p}(\sum_{i= 1}^n b_i^p)^\frac{1}{p}\]
  这里$p > 0 ,q > 0 ,$且$\frac{1}{p}+\frac{1}{q}$

  \[\text{我们取},f(x) = x^{\frac{1}{q}},f''(x) = \frac{1}{q}(\frac{1}{q}-1)x^{\frac{1}{q} - 2} < 0 ,x>0\]
  所以$f(x)$是一个凹函数,若$x_i > 0 ,t_i  > 0 ,\sum_{i = 1}^n t_i = 1$,则
  \[t_i x_1^{\frac{1}{q}} + \cdots + t_n x_n^{\frac{1}{q}}  \leq  (t_1 x_1 + \cdots + t_n x_n )^{\frac{1}{q}}\]

  \[\text{我们取},t_i = \dfrac{a_i^p}{\sum_{i=1}^n a_i^p} ,x_i = \dfrac{b_i^p}{a_i^p} , p - \frac{p}{q} = 1\]
  \[t_i x_i^\frac{1}{q} = {\dfrac{a_i b_i}{\sum_{i = 1}^n a_i^p}}\]
  \[t_i x_i = \dfrac{b_i^q}{\sum_{i = 1}^n a_i^p}\]

  带回原式即可完成证明


\end{example}

\subsubsection{拐点}
凹凸性改变的点

\begin{identification}
$f$在$U(x_0,\delta)$连续且$f(x)在x_0$两侧凹凸性相反,称$x_0$为函数$f$的拐点
\end{identification}
这是一个局部性的概念\\

这只是一个充分条件\\
\begin{theorem}
$x_0$为拐点,且$f''(x_0)$存在,则$f'(x_0) = 0$
\end{theorem}
\begin{cproof}
$f''(x_0)$存在$\Rightarrow f'(x)$在$U(x_0,\delta)$上存在
$x_0$左右两侧凹凸性相反 $\Rightarrow$ 导数左升右降/左降右升\\
$\Rightarrow x_0是f'(x)$的极值点$\Rightarrow f''(x_0) = 0$
\end{cproof}

\begin{theorem}
  \[f''(x_0) = 0 ,f'''(x_0) \text{存在且不为}0,\text{则}x_0\text{为f的拐点}\]
\end{theorem}
\begin{cproof}
\[f'''(x_0) \exists \Rightarrow f''(x)\text{在}U(x_0)\text{中存在}\]
\[f''(x) = f''(x_0) + f'''(x_0)(x-x_0) + o(x-x_0)\]
\[\text{由于}x_0\text{两边}f''(x)\text{异号,从而凹凸性改变,}x_0\text{是拐点}\]
\end{cproof}
\subsubsection{方程求解的牛顿迭代法}
求$f(x),f(c)= 0,c$是解,$f$在$(a,b)$上二阶可导,$f''(x) \neq 0,f'(x) \neq 0$(由Dabourx可知严格单调且凹凸性不变)\\

$f\in C[a,b] ,f(a) f(b) < 0$\\
那么$c$是唯一的解\\

\begin{center}
  \begin{tikzpicture}
    \begin{axis}[
      xlabel=$x$,
      ylabel=$f(x)$,
      axis lines=middle,
      xmin=0, xmax=1.2,
      ymin=-0.5, ymax=2,
      width=\textwidth,
      height=8cm,
      xtick = \empty,
      ytick = \empty
    ]
    \addplot[black,line width=1pt,domain=0.5:1.3,samples=200] {2.5*(x)^2 - 1};
    \addplot[only marks, mark=*] coordinates {(1, {2.5*(1)^2 - 1}) (0.7, 0) (0.7, {2.5*(0.7)^2 - 1}) (0.645,0)};

    \addplot[blue, line width=1pt] coordinates {(1, {2.5*(1)^2 - 1}) (0.7, 0)};
    \addplot[red, line width=1pt] coordinates {(0.7, {2.5*(0.7)^2 - 1}) (0.645, 0)};

    % \addplot[red, line width=1pt] coordinates {(0.6, {tan(deg(0.6))}) (0.9, {tan(deg(0.9))})};
    
    % \addplot[red, line width=1pt] coordinates {(0.9, {tan(deg(0.9))}) (1.2, {tan(deg(1.2))})};

    \node at (axis cs:1, {2.5*(1)^2 - 1}) [anchor=north west] {$(x_1,f(x_1))$};

    \node at (axis cs:0.7, 0) [anchor=north] {$x_2$};
    \node at (axis cs:0.635, 0) [anchor=north] {$x_3$};
    \node at (axis cs:0.7, {2.5*(0.7)^2 - 1}) [anchor=south east] {$(x_2,f(x_2))$};

    \end{axis}
  \end{tikzpicture}
\end{center}
\begin{theorem}
  (1)$f'(x) \text{在}(a,b)\text{上不为}0$\\
  (2)$f(c) = 0 c\in (a,b)$\\
  (3)$x_1\in (a,b) f(x_1)f''(x_1) > 0$\\
  令$x_{n+1} - x_n = \dfrac{f(x_n)}{f'(x_n)}$\\
  \[f'(x) f''(x) > 0 \text{时} x_n\downarrow \to c,f'(x) f''(x) < 0 \text{时} x_n\uparrow \to c\]
\end{theorem}
取决于一阶导和二阶导不能变号,而且$x_1$的选取也有一定要求(满足条件很多,否则我们得到的序列很可能不收敛)\\
牛顿法的收敛速度很快\\
\begin{cproof}
设$f' > 0 ,f'' > 0$,选取$x_1$ ,使得$f(x_1) > 0$,由于单调性$x_1 > c$\\

\[x_2 = x_1 - \dfrac{f(x_1)}{f'(x_1)} < x_1\]

由于$f$是严格凸的
\[0 = f(c) > f(x_1) + f'(x_1) (x_1 - c)\]
从而$x_2 > c$,\\
由归纳法 $c < x_k < x_{k - 1}$ 且有$f(x_k)  > 0$
由于f严格凸,我们同样可以得到
\[0 = f(c) > f(x_k) + f'(x_k) (x_k - c)\]
从而$c < x_{k+1} < x_k$,我们地道道序列单调下降有下界,由单调收敛原理,该序列有极限

我们记\[\lim_{n\to \infty x_n} = \xi\]
\[\xi = \xi - \dfrac{f(\xi)}{f'(\xi)} f(\xi) = 0 , \xi = c\]
其他情况同理可证明\\
\end{cproof}

\section{不定积分}
\subsection{原函数和不定积分的概念}
\begin{identification}
  区间$X$上给定的$f(x),\exists F(x)$满足\[F'(x) = f(x),\forall x\in X, dF(x) = f(x) dx\]\\
  称$F(x)$为$f(x)$的原函数,$f(x)$的所有原函数成为$f$的不定积分,记为$\int f(x)dx$\\

  若$f(x)$有原函数,称$f(x)$是可积的
\end{identification}
[有些时候原函数存在但是写不出来,以及不定积分的可积和定积分的可积不是一个概念]\\

原函数不唯一,$F(x)$为原函数,则$F(x) + C$ 也是原函数

\begin{theorem}
  $F(x)$为$f(x)$的一个原函数,则$\int f(x) dx = F(x) + C,C$为任意常数
  $F(x) + C$ 是$f(x)$的原函数
\end{theorem}

\begin{cproof}
$G(x)$为另一原函数 $G'(x) = f(x), F'(x) = f(x)$\\
\[G'(x) - F'(x) = 0 ,\therefore G(x) - F(x) = C (Rolle \quad Theorem) \]

\end{cproof}
形如$F(x) + C$这样的曲线成为$f(x)$的积分曲线\\

求导是线性的运算-不定积分的线性性质\\
\begin{theorem}
  \[F'(x) = f(x) ,G'(x) = g(x)\]
  \[[\alpha F(x) + \beta G(x)]'  = \alpha F'(x) + \beta G'(x) = \alpha f(x) + \beta g(x)\]

  从而我们得出
  \[\int(\alpha f(x) + \beta g(x))dx = \alpha F'(x) + \beta G'(x) + C = \alpha \int f(x)dx + \beta \int g(x)dx\]
\end{theorem}
这也是其唯一的规律\\

\begin{example}
  \[\int \dfrac{dx}{\sin^2 2x} = \int\frac{1}{4}\dfrac{\sin^2 x+\cos^2 x}{\sin^2 x \cos^2 x} = \frac{1}{4}\int (\dfrac{1}{\sin^2 x} + \dfrac{1}{\cos^2 x} ) = \frac{1}{4}(\tan x - \cot x) + C \]
\end{example}

\begin{example}
  \[\int \frac{1}{x}dx = \ln|x| + C\]
\end{example}


\subsection{换元积分法}
\paragraph{一阶微分的形式不变性}
如果$u$是自变量
\[dF(u) =  F'(u)du\]
如果$u$是中间变量
\[dF(u(x)) = F'(u(x)) u'(x)dx = F'(u(x))du\]

\subsubsection{第一换元法}
把$d$看作微分符号\\
\begin{theorem}
  如果$u$是自变量
  \[\int f(u)du = F(u) + C\]
  如果$u$是中间变量
  \[ \int f(u(x)) u'(x)dx = F(u(x))  + C\]
\end{theorem}

\begin{example}
  \[\int \dfrac{dx}{x-1} = \int \dfrac{d(x-1)}{x-1}=\int \dfrac{d(u)}{u}=\ln|x-1| + C \]
\end{example}

\begin{example}
  \[\int \dfrac{xdx}{1+x^2} = \dfrac{1}{2}\int \dfrac{dx^2}{1+x^2}=\dfrac{1}{2}\int \dfrac{d(x^2+1)}{1+x^2}= \dfrac{1}{2}\ln|x^2 + 1 | + C\]
\end{example}

\begin{example}
  \[\int \dfrac{dx}{a^2 + x^2} =\dfrac{1}{a}\int \dfrac{d(\frac{x}{a})}{1 + (\frac{x}{a})^2} = \dfrac{1}{a} \arctan (\dfrac{x}{a}) + C\]
\end{example}

\begin{example}
  \[\int \dfrac{dx}{\sqrt{a^2 - x^2}} = \int \dfrac{d(\frac{x}{a})}{\sqrt{1-(\frac{x}{a})^2}} = \arcsin (\dfrac{x}{a}) + C\]
\end{example}

\paragraph{一些三角函数的例子}
\begin{example}
  \[\int \sin^3 x = \int \sin^2 x \sin x dx = -\int(1- \cos^2) d(\cos x)\]
\end{example}

\begin{example}
  \[\int \sin^2 x dx = \int \dfrac{1-\cos 2x}{2} dx = \dfrac{1}{x} - \dfrac{1}{4}\int \cos 2x\, d(2x)\]
\end{example}
\begin{example}
  \[\int \tan x dx = \int \dfrac{\sin x}{\cos x} dx = \int \dfrac{d(\cos x )}{\cos x}\]
\end{example}

\begin{example}
  求$I_n = \int \tan^n x dx$\\
  \[
  \begin{aligned}  
  I_n &= \int \tan^{n-2} \dfrac{\sin^2 x}{\cos^2 x}dx = \int \tan^{n-2} [\dfrac{1}{\cos^2 x} - 1]dx \\
  &= \int \tan^{n-2} d(\tan x) - \int \tan^{n-2}dx\\ 
  &= \dfrac{1}{n-1} \tan^{n-1} x - I_{n-2}\\
  \end{aligned}
  \]
  我们就找到了这样的递推式,只要找到前两项即可\\
\end{example}

\begin{example}
  \[\int \dfrac{dx}{\cos^2 x+ 1} = \int \dfrac{dx}{(\sec^2 x+ 1)\cos^2 x} = \int \dfrac{d(\tan x)}{2 + \tan^2 x}\]
\end{example}

\paragraph{构造方程求解不定积分}

\begin{example}
  求$\int \dfrac{dx}{1 + x^3}$和$\int \dfrac{xdx}{1 + x^3}$\\
  \[
  \begin{aligned}  
  F(x) + G(x) &= \int \dfrac{1 + x}{1 + x^3}dx = \int \dfrac{1}{1-x + x^2}dx\\
  &=\int \dfrac{d(x-\frac{1}{2})}{\frac{3}{4}+(x-\frac{1}{2})^2}\\
  &=\int \dfrac{\frac{4}{3}d(x-\frac{1}{2})}{1+[\frac{2}{\sqrt{3}}(x-\frac{1}{2})]^2}\\
  &=\dfrac{2}{\sqrt{3}}\int \dfrac{d(\frac{2}{\sqrt{3}}(x-\frac{1}{2}))}{1+[\frac{2}{\sqrt{3}}(x-\frac{1}{2})]^2}\\
  &=\dfrac{2}{\sqrt{3}}\tan\left(\dfrac{2x - 1}{\sqrt{3}}\right)\\
  &\\
F(x) - G(x) &= \int \dfrac{1 - x}{1 + x^3}dx = \int \dfrac{1-x + x^2 - x^2}{(1+ x )(1 -x + x^2)}dx \\
&= \int \dfrac{dx}{1+x} - \int \dfrac{x^2}{1+x^3} dx  \\
&= \int \dfrac{dx}{1+x} - \dfrac{1}{3}\int \dfrac{dx^3}{1+x^3} \\
&=\ln|1 + x| - \dfrac{1}{3}\ln|1 + x^3|\\
  \end{aligned}
\]
\end{example}
\subsubsection{第二换元法}
\begin{theorem}
设函数$f(x)$在某一开区间上可导,$x'(t) \neq 0$,如果
\[\int f(x(t))x'(t)dt = G(t) + C\]
那么
\[\int f(x)dx = G(t(x)) + C\]
其中$t(x)$为$x(t)$的反函数\\
\end{theorem}

第二换元法主要用来求解含有根式的积分\\
\begin{example}
\[\int \sqrt{a^2 - x^2}\]
令 $x = a\sin t, t\in [-\dfrac{\pi}{2},\dfrac{\pi}{2}]$,则$dx = a \,d(\sin t)= a\cos t \, dt$
\[\int \sqrt{a^2 - x^2} = \int a\cos t \,d( a\sin t) = \int a^2 \cos^2 t dt = a^2 \int \dfrac{1 + \cos 2t}{2} dt\]
\end{example}

\begin{example}
  \[\int \dfrac{dx}{\sqrt{x^2 - a^2}}\]
  令 $x = a\sec t, t\in(-\dfrac{\pi}{2},\dfrac{\pi}{2})$,则$dx = a \,d(\sec t) = a~ \dfrac{\sin t}{\cos^2 t}\, dt = a \sec t \tan t \, dt$
  \[\int \dfrac{dx}{\sqrt{x^2 - a^2}} =\int \dfrac{a\sec t \tan t \, dt}{a \tan t} = \int \dfrac{dt}{\cos t}\]
  \[
  \begin{aligned}  
  \int \dfrac{dt}{\cos t} &= \int \dfrac{\cos t}{\cos^2 t}\, dt \\
  &= \int \dfrac{d(\sin t)}{1- \sin^2 t} \\
  &=\dfrac{1}{2}\ln \left| \dfrac{1 + \sin t}{1 - \sin t}\right| + C\\ 
  &=\dfrac{1}{2}\ln \left| \dfrac{(1 + \sin t)(1 + \sin t)}{(1 - \sin t)(1 + \sin t)}\right| + C\\ 
  &=\dfrac{1}{2}\ln \left| \dfrac{1 + \sin t}{\cos t}\right|^2 + C\\ 
  &=\ln \left|\sec t + \cot t \right| + C\\ 
  &= \ln| x + \sqrt{x^2 - a^2}| + C\\
  \end{aligned}
  \]
\end{example}

\begin{example}
  \[\int \dfrac{dx}{\sqrt{x^2 + a^2}}\]
  令 $x = a\tan t, t\in(-\dfrac{\pi}{2},\dfrac{\pi}{2})$,则$dx = a ~\dfrac{1}{\cos^2 t}dt$
  \[\int \dfrac{dx}{\sqrt{x^2 + a^2}} =\int \dfrac{dt}{\sec t \cos ^2 t }= \int \dfrac{dt}{\cos t} = \ln \left| \dfrac{1 + \sin t}{\cos t}\right| + C = \ln\left| x + \sqrt{x^2 + a^2}\right| + C\]
\end{example}

\begin{example}
  
\end{example}

\subsection{分部积分}

\begin{theorem}
  \[u(x),v(x)\text{可导},\int u'(x)v(x) dx\text{存在}\]
  \[then,\int u(x)v'(x) dx= u(x)v(x) - \int u'(x) v(x)dx\]
  \[\int u dv = uv- \int v du\]
\end{theorem}

\begin{cproof}
\[
\begin{aligned}  
&\because (uv)' = u'v + uv'\\
&\therefore\int(uv)'dx = \int u' v dx + \int u v' dx\\
&\therefore uv = \int u' v dx + \int u v' dx\\
&\therefore \int u dv = uv- \int v du\\
\end{aligned}
\]
\end{cproof}
分部积分主要处理的是乘积形式的积分,目的是让求导越求越简单\\
\begin{example}
  \[
  \begin{aligned}  
  \int x^2 \sin x \, dx &= -\int x^2 \cos x\, dx = -x^2 \cos x + \int \cos x d(x^2) \\
  &= -x^2 \cos x + 2\int x\cos x dx \\ 
  &= -x^2 \cos x + 2\int x d(\sin x)\\
  &= -x^2 \cos x + 2 x \sin x - 2 \int \sin x\, dx  \\
  &=-x^2 \cos x + 2 x \sin x + 2 \cos x + c\\
  \end{aligned}
  \]
\end{example}

\begin{example}
  \[
  \begin{aligned}  
  \int \sqrt{x^2 - 1}\, dx &= x\sqrt{x^2 -1} - \int x d(\sqrt{x^2 -1})\\
  &= x\sqrt{x^2 -1} - \int \dfrac{x^2 dx}{\sqrt{x^2 -1}} \\
  &= x\sqrt{x^2 -1} - \int\sqrt{x^2 -1}\, dx - \int \dfrac{dx}{\sqrt{x^2 -1}}\\
  \end{aligned}
  \]
  我们移项可得
  \[
    \begin{aligned}  
    2\int \sqrt{x^2 - 1} \,dx &= x\sqrt{x^2 -1} - \int \dfrac{dx}{\sqrt{x^2 -1}}\\
    &= x\sqrt{x^2 -1} - \ln\left|x + \sqrt{x^2 - 1} \right|
    \end{aligned}
    \]
  \[
    \begin{aligned}  
    \int \sqrt{x^2 - 1} \,dx &= \dfrac{1}{2}x\sqrt{x^2 -1} - \dfrac{1}{2}\ln\left|x + \sqrt{x^2 - 1} \right|
    \end{aligned}
    \]

\end{example}

\begin{example}
  \[
  \begin{aligned}  
  \int e^{ax} \cos bx dx &= \int \dfrac{1}{a} \cos bx\, d(e^{ax}) \\
  &= \dfrac{1}{a} \cos bx e^{ax} - \dfrac{1}{a}\int e^{ax} d(\cos bx)  \\
  &= \dfrac{1}{a} \cos bx e^{ax} + \dfrac{b}{a}\int e^{ax} \sin bx \,dx\\
  &= \dfrac{1}{a} \cos bx e^{ax} + \dfrac{b}{a^2}\int  \sin bx \,d(e^{ax})\\
  &= \dfrac{1}{a} \cos bx e^{ax} + \dfrac{b}{a^2}\sin bx~ e^{ax} - \int \dfrac{b}{a^2} e^{ax} \,d(\sin bx)\\
  &= \dfrac{1}{a} \cos bx e^{ax} + \dfrac{b}{a^2}\sin bx~ e^{ax} - \int \dfrac{b^2}{a^2} e^{ax}\cos bx \,dx\\
  \end{aligned}
  \]
  之后移项解方程即可\\
\end{example}

\subsection{有理函数的积分}
有理函数$\dfrac{p(x)}{q(x)} = \text{多项式} + \text{真分式}$\\

真分式 $\Rightarrow$ 简单分式之和\\
\[\dfrac{A}{(x-a)},\dfrac{A}{(x-a)^m} \text{有实根}\]
\[\dfrac{Bx + c}{x^2 + px + q},p^2 - 4q < 0,\dfrac{Bx + c}{(x^2 + px + q)^m},\text{有虚根}\]

怎么拆?\\
假设$p(x)$是$(n-1)$阶多项式,$q(x)$是$n$阶多项式,那么这个时候"多项式"部分不存在,即仅存在"真分式"\\

\[q(x) = \prod_{i = 1}^r (x- a_i)^{h_i}\prod_{j = 1}^s (x^2 + px + q)^{k_j},\text{其中} \sum_{i = 1}^r h_i+ \sum_{j = 1}^s 2k_j = n\]

$p(x)$的系数共有$n$个(在$x^{n-1},\cdots, x^0$前面)\\

\begin{theorem}
上述$\dfrac{p(x)}{q(x)}$形式的真分式,可以唯一的表示为如下的简单分式之和(证明见高等代数书籍)
\[
\begin{aligned}  
\dfrac{p(x)}{q(x)} &= \sum_{i = 1}^r \left(\dfrac{A_i}{(x-a_i)} + \dfrac{A_i^{1}}{(x-a_i)^{2}} + \cdots + \dfrac{A_i^{h_i-1}}{(x-a_i)^{h_i}}\right) \\
&+ \sum_{j = 1}^s\left(\dfrac{M_j x + N_j }{(x_j^2 + p_j x + q_j)} + \dfrac{M_j^1 x + N_j^1 }{(x^2 + p_j x + q_j)^2} + \cdots +\dfrac{M_j^{(k_j -1)} x + N_j^{(k_j -1)} }{(x^2 + p_j  x + q_j)^{k_j}}\right)
\end{aligned}
\]
\end{theorem}

其中的未知数为\[\sum_{i = 1}^r h_i+ \sum_{j = 1}^s 2k_j =n\]个,因此是可解的\\

\begin{example}
  \[\dfrac{x+3}{(x + 2 )(x^2 - 1 )} = \dfrac{A}{x + 2}+\dfrac{B}{x + 1}+\dfrac{C}{x -1}\]
  方法一:通分后解方程组\\

  方法二:求极限\\
  两边乘上$(x + 2)$,再令$x \to -2$,从而解出$A = \dfrac{1}{3}$\\
  两边乘上$(x + 1)$,再令$x \to -1$,从而解出$B = -1$\\
  同样求得C
\end{example}
有重根的例子
\begin{example}
  \[\dfrac{x + 3}{(x + 1) (x - 2)^2} = \dfrac{A}{x + 1} + \dfrac{B}{x-2} + \dfrac{C}{(x-2)^2}\]

求A的方法可以仿照上例获得

两边同时乘上$(x-2)^2$可以求得C

与原先唯一的区别是,我们怎么求B呢?主要有以下2个方法\\

(1)两边同时乘以$x$,那么
\[\dfrac{(x + 3)x}{(x + 1) (x - 2)^2}=\dfrac{A}{x + 1}x + \dfrac{B}{x-2}x + \dfrac{C}{(x-2)^2}x\]
再令$x\to +\infty$
\[\lim_{x\to +\infty}\dfrac{(x + 3)x}{(x + 1) (x - 2)^2}=\lim_{x\to +\infty}\left(\dfrac{A}{x + 1}x + \dfrac{B}{x-2}x + \dfrac{C}{(x-2)^2}x\right)\]
得到
\[ 0 = A + B \]
于是我们就求得了B\\

(2)两边乘上$(x-2)^2$后
我们得到
\[\dfrac{x + 3}{(x + 1)} = \dfrac{A(x - 2)^2}{x + 1} + B(x-2) + C\]
对$x$求导
\[\dfrac{-2}{(x + 1)^2} = \dfrac{A(x^2 + 2x -8)}{(x + 1)^2} + B \]
再令$x = 2$\\
即可解得$B = -\dfrac{2}{9}$
\end{example}
如果出现复根(重根)
\begin{example}
\[\dfrac{2x+2}{(x-1)(x^2 + 1)^2} = \dfrac{A}{x-1} + \dfrac{bx + C}{(x^2 + 1)} +  \dfrac{Dx + E}{(x^2 + 1 )^2} \]
首先两边同时乘上$x-1$,之后令$x\to 1$,可以解得$A$\\

然后两边同时乘以$(x^2 + 1)^2$,我们得到
\[\dfrac{2x+2}{x+1} = (Bx + C)(x^2 + 1) + Dx + E\]
令$x\to i$(虚数单位)\\
\[Di + E = \dfrac{2i + 2}{i - 1}\]
从而可以解出$D,E$\\

最后,两边同乘$x$,令$x\to + \infty$(这个操作非常常用),得到\\
\[A+B = 0\]
同样可以令$x = 0$,那么可得$C$\\
\end{example}
\subsubsection{简化分式的不定积分}

我们前面已经学会了如何将有理分式化简为简化分式,之后我们要考虑如何实现积分的操作\\
\begin{equation*}
  \begin{cases}
    \int\dfrac{A}{(x-a)}\,dx = \ln|x-a| + C\\
    \int\dfrac{A}{(x-a)^m}\,dx = \dfrac{-1}{m-1}\dfrac{A}{(x-a)^{m-1}} + C\\
  \end{cases}
\end{equation*}
接下来我们考虑的情况
\[\dfrac{Bx + c}{x^2 + px + q},p^2 - 4q < 0,\dfrac{Bx + c}{(x^2 + px + q)^m},\text{有虚根}\]

\[
\begin{aligned}  
\int\dfrac{Bx + C}{(x^2 + px + q)^k}&=\int\dfrac{Bx + C}{[(x+\frac{p}{2})^2 + q - \frac{p^2}{4}]^k}\\
&=\int\dfrac{Bx + C - \frac{Bp}{2}}{(t^2 + a^2)^k}\text{进行变量代换}((x+\frac{p}{2})^2 = t^2 ,q - \frac{p^2}{4} = a^2 )\\
&=B\int \dfrac{Bx + C}{(t^2 + a^2)^k}\, dt + \int(C - Bp)\dfrac{dt}{(t^2 + a^2)^k}\\
&k = 1\text{时}\\
&\int \dfrac{Bx + C}{(t^2 + a^2)^k} = \dfrac{1}{x}\ln|t^2 + a^2 | + C\\
&\int \dfrac{dt}{(t^2 + a^2)^k}=\dfrac{1}{a}\arctan \dfrac{t}{a} + C\\
&k \neq 1\text{时我们有,上面形式的式子积分可以直接得到}\\
&\text{而下面形式的式子可以构造递推公式得到,在此我们省略}\\
\end{aligned}
\]
\subsubsection{三角函数有理式}
\[R(\sin x,\cos x)\]
其中\[R(\sin x,\cos x)\]是有理函数,即\[R(\sin x,\cos x) =\dfrac{P(u,v)}{Q(u,v)}\]
其中$P,Q$为多项式,且每一项的形式$Cu^i v^j,i,j\in \NN$\\
\paragraph{万能代换}
我们做变量代换$\tan \frac{x}{2} = t$从而
\[\sin x = \dfrac{2t}{1 + t^2},\cos x =  \dfrac{1-t^2}{1+t^2}\]
最终我们得到之前的有理分式的形式,因此三角函数有理式都是可积的\\

然而,很多时候,一些三角函数有理式并不需要这样的换元可能更加方便,例如

当$R(\sin x,\cos x)$关于$\sin x$是奇函数时,我们采用$t = \cos x$\\

当$R(\sin x,\cos x)$关于$\cos x$是奇函数时,我们采用$t = \sin x$\\

当$R(\sin x,\cos x)$关于$\sin x,\cos x$都是偶函数时,我们采用$t = \tan x$\\

\begin{example}
  \[
    \dfrac{dx}{a^2 \sin^2 x + b^2 \cos^2 x}=\dfrac{d(\tan x)}{a^2 \tan^2 x + b^2}=\dfrac{dt}{a^2 t^2 + b^2}(\text{令} t = \tan x)\\ 
  \]
\end{example}

\section{定积分}
\subsection{定义与基本性质}
\begin{identification}
$f$在$[a,b]$上定义,记$P$是$[a,b]$上一分割
\[P: a = x_0< x_1 <x_2\cdots < x_n = b\]
$x_i$称为分割点,我们定义
\[\Delta x_i = x_i - x_{i-1}, i =1,2,\cdots n\]
定义
\[|P| = \max_{i = 1,2 \cdots, n}\{\Delta x_i\}\]
称为分割宽度,而
$\xi_i \in [x_{i-1},x_i]$称为标志点
而\[\sum_{i = 1}^n f(\xi_i)\Delta x_i = \sigma(f,P,\xi)\]
称为黎曼和\\
\end{identification}

我们要考虑的是极限,即
\[\lim_{|P|\to 0} \sigma(f,P,\xi) \equiv I := \int_{a}^{b}f(x)\,dx\]

如果极限$I$存在,则称$f$在$[a,b]$上黎曼可积,$I$为$f$在$[a,b]$上定积分\\

定义条件:

(1)$[a,b]$是有界闭区间(尽管说整个实轴既是开区间也是闭区间,但不是有界的)

(2)极限与划分$P$的方式无关,与标志点$\xi$的选取无关,只与划分的宽度$|P|$有关

\begin{identification}$(\varepsilon-\delta)$
  \[\forall \varepsilon > 0,\exists \delta > 0 ,\forall |P| < \delta ,\forall \xi ,s.t. |\sigma(f,P,\xi) - I|< \varepsilon \]

  其中
  \[I = \lim_{|P|\to 0} \sigma(f,P,\xi) \equiv I := \int_{a}^{b}f(x)\,dx\]
\end{identification}

连续函数是黎曼可积的\\

\subsubsection{定积分的性质}

\begin{theorem}(①线性)
如果$f(x),g(x)$在$[a,b]$上黎曼可积,那么
\[\forall \alpha,\beta\in \RR , \alpha f(x) + \beta g(x)\text{在}[a,b]\text{上黎曼可积}\]
\[\text{且} \int_{a}^{b} (\alpha f(x) + \beta g(x))\,dx = \alpha \int_{a}^{b}f(x)\,dx + \beta \int_{a}^{b}g(x)\,dx\]
\end{theorem}

\begin{lemma}
  $f$在$[a,b]$上可积,则$f$在$[a,b]$上有界\\
\end{lemma}

\begin{cproof}
  反证,任意分割$P: a = x_0< x_1 <x_2\cdots < x_n = b$\\
  假设$f$无界,那么至少存在某个$[x_{k-1},x_k]$
  \[\sum_{i = 1}^n f(\xi_i)\Delta x_i  = f(\xi_k)\Delta x_k+\sum_{i \neq k}^n f(\xi_i)\Delta x_i \]

  取定$x_l,l\neq k$

  取$\xi_k$,~ s.t.$| f(\xi_k)\Delta x_k|$可以大于任何的数

  \[\Rightarrow \sum_{i = 1}^n f(\xi_i)\Delta x_i\text{可以大于任何的数}\]

  因此这个黎曼和没有极限,与不可积矛盾
\end{cproof}

\begin{example}
  \begin{equation*}
    f(x)=
    \begin{cases}
      1,x\text{为无理数}\\
      0,x\text{为有理数}
    \end{cases}
    x\in[0,1]
  \end{equation*}

显然$f$有界

然而,无论怎么分割,我们都可以取$\xi_i$全为无理数或有理数,前者得到黎曼和为$1$,而后者得到黎曼和为$0$,在这种情况下黎曼和与$\xi$的选取有关,因此这个函数不是黎曼可积的\\
\end{example}

\begin{theorem}(②积分可加性)\\

  $f$在$[a,b]$上可积,在$[b,c]$上可积$\Rightarrow f$在$[a,c]$上可积,且
  \[\int_{a}^{c}f(x)\,dx=\int_{a}^{b}f(x)\,dx+\int_{b}^{c}f(x)\,dx\]
  
\end{theorem}

\begin{cproof}
$f$在$[a,b]$,$[b,c]$上有界\\
\[|f(x)|<k,\forall x\in[a,c]\]

我们作$[a,c]$上任意一个分割$P: a = x_0< x_1 <x_2\cdots < x_n = c$

我们要考虑两种情况

①$b$是$P$的一个分点,则取$\tilde{P} = P ,\tilde{\xi} = \xi$

②$b$不是$P$的一个分点,这个时候我们要添加$b$为一个分点,同时也把b作为左右两个分割区间的标志点,即有\\

\[\tilde{P}: a = x_0< x_1 <\cdots<x_{k-1} < b< x_k<\cdots < x_n = c\]
\[\tilde{\xi} = (\xi_1,\cdots,\xi_{k-1},b,b,\xi_{k+1},\cdots,\xi_n)\]
$[a,c]$上的黎曼和
\[\sigma(f,P,\xi) \rightarrow f(\xi_k)(x_k - x_{k-1})\]
而现在是
\[\sigma(f,\tilde{P},\tilde{\xi})\rightarrow f(b)(x_k - x_{k-1})\]
差异在$[x_{k-1},x_{k+1}]$这一段\\

二者相减
\[\left|f(\xi_k)(x_k - x_{k-1})-f(b)(x_k - x_{k-1})\right| = \left|f(\xi_k)-f(b)\right|(x_k - x_{k-1}) < 2k |P|\]

\[\sigma(f,\tilde{P},\tilde{\xi}) = \sigma(f,\tilde{P'},\tilde{\xi'}) + \sigma(f,\tilde{P''},\tilde{\xi''})\]

\[\lim_{|P|\to 0}\left|f(\xi_k)(x_k - x_{k-1})-f(b)(x_k - x_{k-1})\right| = 0\]
由前面的
\[\sigma(f,\tilde{P},\tilde{\xi}) = \sigma(f,\tilde{P'},\tilde{\xi'}) + \sigma(f,\tilde{P''},\tilde{\xi''})\]
我们有
\[\lim_{|P|\to 0}|\sigma(f,P,\xi) =\int_{a}^{c}f(x)\,dx=\int_{a}^{b}f(x)\,dx+\int_{b}^{c}f(x)\,dx\]
\end{cproof}
与此同时,我们约定
\[\int_{a}^{c}f(x)\,dx=-\int_{c}^{a}f(x)\,dx\]

事实上,对于任意的$a,b,c\in \RR$,如果$f(x)$在这较大的区间中可积,那么
\[\int_{a}^{b}f(x)\,dx + \int_{b}^{c}f(x)\,dx=\int_{a}^{c}f(x)\,dx\]

\begin{theorem}
  (积分的单调性)

  如果函数$f(x),g(x)$在区间$[a,b]$上黎曼可积,并且$f(x) \leq g(x)$,则
  \[\int_{a}^{b}f(x)\,dx\leq \int_{a}^{b}g(x)\,dx\]
\end{theorem}

\begin{theorem}
  (积分的中值定理)

  如果函数$f(x)$在$[a,b]$上可积,那么($f(x)$在$[a,b]$上有界)
  \[\exists m,M \in \RR,  m \leq f(x) \leq M\]
  从而
  \[m(b - a) \leq\int_{a}^{b}f(x)\,dx \leq M (b - a)\]
  特别的,如果$f(x)\in C[a,b]$,那么
  \[\exists\xi \in [a,b],s.t.f(\xi) = \dfrac{\int_{a}^{b}f(x)\, dx}{b - a}\]
\end{theorem}
\subsection{牛顿-莱布尼茨公式}
\begin{theorem}
  (牛顿-莱布尼茨公式)
  如果$f\in C[a,b]$,且存在函数函数$F(x) \in [a,b]$,且在$(a,b)$上可导,使得函数\[F'(x) = f(x)\]
  那么
  \[\int_{a}^{b}f(x)\,dx = F(b) - F(a)\]

\end{theorem}

\begin{example}
  \[\lim_{n\to +\infty} \dfrac{1^p + 2 ^ p +\cdots + n^p}{n^{p + 1}}\]

  可以写成黎曼和的形式
  \[
  \begin{aligned}  
  \lim_{n\to +\infty}\sum_{k = 1}^{n} \dfrac{k^p}{n^p} \cdot \dfrac{1}{n}&=\lim_{n\to +\infty}\sum_{k = 1}^{n} \xi_k^p \Delta x_k\\
  &=\lim_{n\to +\infty}\sum_{k = 1}^{n} \xi_k^p \Delta x_k\\
  &=\int_{0}^{1}x^p\,dx 
  \end{aligned}
  \]
  
\end{example}


\begin{identification}
连续可导(可微)

$\varphi(t)$在$(a,b)$上每一点可导,$\varphi'(t) \in C(a,b)$,则称$\varphi(t)$在$(a,b)$上连续可导,记作
\[\varphi(t) \in C^1 (a,b)\]

对于闭区间,如果$\varphi(t) \in C^1 (a,b)$
且有
$\varphi(t)\text{在}a\text{右侧可导},\text{在}b\text{左侧可导,且} \varphi'(t) \in C[a,b]$则称$\varphi(t)$在$[a,b]$上连续可导,记作
\[\varphi(t) \in C^1 [a,b]\]
\end{identification}

\begin{theorem}
  (定积分的换元法)

  设$\varphi(t) \in C^1[\alpha,\beta]$,且有$\varphi(\alpha) = a , \varphi(\beta) = \beta , \varphi((\alpha,\beta)) \subset (a,b),f(x) \in C[a,b]$(这里没有要求$\varphi'(t) \neq 0$)则
  \[\int_{a}^{b}f(x)\,dx = \int_{\alpha}^{\beta}f(\varphi(t))\varphi'(t)\,dt\]

\end{theorem}
为什么不用要求$\varphi'(t) \neq 0$?

因为在定积分中,我们不用把$t$代换为$x$,只需要在开始时将$x$的范围映射到$t$的范围,然后对$t$求积分即可得到数值\\

定积分有没有"第一换元法"?事实上,这样的操作实际是通过不定积分的第一换元法求出原函数,再利用牛顿-莱布尼茨公式求出积分,而这里的[定积分的换元法]是不同于不定积分换元法的方式,因为这里改变了积分的上下限\\
\begin{cproof}
$f\in C[a,b] \Rightarrow f$有原函数记为$F(x)$,使得

\[F'(x) = f(x)\]
根据求导法则
\[\dfrac{d(F(\varphi(t)))}{dt} = F'(x)\varphi'(t) = f(\varphi(t))\varphi'(t)\]
\[
\begin{aligned}
\int_{\alpha}^{\beta}f(\varphi(t))\varphi'(t)\,dt &= F(\varphi(t))|_{\alpha}^{\beta}\\
&=F(b) - F(a)\\
&=\int_{a}^{b}f(x)\,dx 
\end{aligned}
\]
\end{cproof}

\begin{example}
  \[
  \begin{aligned}
  \int_{0}^{1}\sqrt{1 - x^2}\, dx &= \int_{0}^{\frac{\pi}{2}} \cos^2 t \, dt\\
  &=\int_{0}^{\frac{\pi}{2}} \dfrac{1 + \cos 2t}{2} \, dt\\
  &=\left(\dfrac{1}{2}t + \dfrac{1}{4}\sin 2t\right)\Big\vert_0^{\frac{\pi}{2}}\\
  &=\dfrac{\pi}{4}\\
  \end{aligned}
  \]
\end{example}
\begin{theorem}
  (分部积分)
  \[\int_{a}^{b}u(x)v'(x)\,dx = u(x)v(x)|_a ^b-\int_{a}^{b}u'(x)v(x)\,dx \]
\end{theorem}

\begin{lemma}
  $\varPsi(t) \in C^{n + 1}[0,1]$,那么
  \[\varPsi(1) = \varPsi(0) + \dfrac{\varPsi'(0)}{1!} +\dfrac{\varPsi''(0)}{2!}+\cdots +\dfrac{\varPsi^{(n)}(0)}{n!} + \dfrac{1}{n!}\int_{0}^{1}\varPsi^{(n + 1)}(0)\, d{(1 - t)^{n + 1}}\]

\end{lemma}

\begin{cproof}
  \[
  \begin{aligned}
\varPsi(1)&= \varPhi(0) + \int_{0}^{1}\varPsi'(t)\,dt\\
&=\varPsi(0) + \int_{0}^{1}\varPsi'(t)\,d(1-t)\\
&=\varPsi(0) -( \varPsi'(t)(1-t)|_0^1 - \int_{0}^{1}(1-t)\,d(\varPsi'(t)))\\
&=\varPsi(0) + \varPsi'(0) - \int_{0}^{1}(1-t)\,d(\varPsi'(t))\\
&=\varPsi(0) + \varPsi'(0) - \int_{0}^{1}(1-t)\,\varPsi''(t)dt\\
&=\varPsi(0) + \varPsi'(0) + \int_{0}^{1}(1-t)\,\varPsi''(t)d(1-t)\\
&\text{重复上述过程(也就是把最右边那项反复进行上述操作),由归纳法可以得到}\\
&= \varPsi(0) + \dfrac{\varPsi'(0)}{1!} +\dfrac{\varPsi''(0)}{2!}+\cdots +\dfrac{\varPsi^{(n)}(0)}{n!} + \dfrac{1}{n!}\int_{0}^{1}\varPsi^{(n + 1)}(0)\, d{(1 - t)^{n + 1}}\\
\end{aligned}
\]
\end{cproof}

\subsubsection{带积分余项的Taylor展开(有限增量)}
\begin{theorem}
如果函数$f(x)$在$x_0$处$n + 1$阶可导,那么
% \[f(x) = f(x_0) + \dfrac{f'(x_0)}{1!}(x - x_0) +\cdots + \dfrac{f^{(n)}(x_0)}{n!}(x - x_0)^n + \dfrac{1}{n!}\int_{0}^{1}(1-t)^n f^{(n + 1)}(x_0 - t(x - x_0))\,dt\]
% 进行变量代换可以得到
\[f(x) = f(x_0) + \dfrac{f'(x_0)}{1!}(x - x_0) +\cdots + \dfrac{f^{(n)}(x_0)}{n!}(x - x_0)^n + \dfrac{1}{n!}\int_{x_0}^{x} f^{(n + 1)}(t) (x - t)^n\,dt\]

由积分余项的泰勒公式进行适当转换可以得到剩下的三种余项
\paragraph{导出拉格朗日余项\\}

对$ \dfrac{1}{n!}\int_{x_0}^{x} f^{(n + 1)}(t) (x - t)^n\,dt$使用第一积分中值定理的变式,记$g(t ) = (x - t)^n$,则$g(t)$在$[x_0,x]$上符号不变,因而$\exists \xi \in (x_0,x)$使得
\[ \dfrac{1}{n!}\int_{x_0}^{x} f^{(n + 1)}(t) (x - t)^n\,dt = \dfrac{f^{(n + 1)}(\xi)}{(n + 1)!}(x - x_0)^{n + 1}\]

\paragraph{导出柯西余项\\}

对$ \dfrac{1}{n!}\int_{x_0}^{x} f^{(n + 1)}(t) (x - t)^n\,dt$使用第一积分中值定理,$\exists \xi \in [x_0,x]$使得
\[\dfrac{1}{n!}\int_{x_0}^{x} f^{(n + 1)}(t) (x - t)^n\,dt = \dfrac{1}{n!}f^{(n + 1)}(\xi) (x - \xi)^n (x - x_0)\]

其中$\xi$可以表述为$\xi = x_0 + \theta (x - x_0),0\leq \theta \leq 1$,带回上式可以得到
\[\dfrac{1}{n!} f^{(n + 1)}(x_0 + \theta(x - x_0))(1 - \theta)^n (x - x_0)^{n + 1}\]
% 如果把$t$换成$\theta$,可以把最后的那一项
% \[\dfrac{(x-x_0)^{n+1}}{n!}\int_{0}^{1}(1-\theta)^n f^{(n + 1)}(x_0 - \theta(x - x_0))\,d\theta\]
% 称为柯西余项\\
\end{theorem}

\begin{cproof}
\[f(x) - f(a) = \int_{a}^{x}f'(t)\,dt\]


将这样的过程推广,我们可以把$0$换为$a$,把$1$换为$x$,重复上述的归纳,
% 得到
% \[f(x) = f(a) + \dfrac{f'(a)}{1!}(x - a) +\dfrac{f''(a)}{2!}(x-a)^2+\cdots + \dfrac{f^{(n)}(a)}{n!} + \dfrac{(x-a)^{n+1}}{n!}\int_{0}^{1}(1-t)^n f^{(n + 1)}(a - t(x - a))\,dt\]
\end{cproof}
\subsection{定积分的几何与物理应用,微元法}
\subsubsection{面积的求解}
正常的积分:可从$x$轴方向也可以从$y$轴方向\\

极坐标下面积积分
\[S = \int_{\alpha}^{\beta}\dfrac{1}{2} r^2(\theta)\,d\theta\]

\subsubsection{体积}
面积关于$x$的积分
\[ V = \int_{a}^{b}A(x)\, dx\]

\begin{example}
  一个函数曲线绕$x$轴旋转得到的旋转体,每个截面都是圆\\
  \[V = \pi \int_{a}^{b}  [f(x)]^2 \, dx\]
\end{example}

\subsubsection{曲线的弧长}

\[L =\int_{\alpha}^{\beta} \sqrt{x'(t)^2 + y'(t)^2}\]

在平面$\RR^2$中,
\begin{equation*}
  \begin{cases}
    x = x(t)\\
    y = y(t)\\ 
  \end{cases}
  \alpha\leq t\leq\beta,x(t),y(t) \in C^1[\alpha,\beta]\text{(或者分段$C^1$)} \\
\end{equation*}
这是为了保证曲线长度有限\\
\[\alpha = t_0 < t_1 < \cdots < t_n = \beta\]
\[\Delta t_i = t_i - t_{i - 1}\]
\[|P_n| = \max\{t_i\}\]

\[\sum_i^n\sqrt{[x(t_i) - x(t_{i- 1})]^2 + [y(t_i) - y(t_{i - 1})]^2} \]
利用拉格朗日定理
\[=\sum_i^n\sqrt{x'(\tau_i)^2 + y'(\tau_i')^2}\Delta t_i \]

我们考察积分和,设其为\[\sum_i^n\sqrt{x'(\tau_i)^2 + y'(\tau_i)^2 }\Delta t_i\]
\[\sum_i^n\sqrt{x'(\tau_i)^2 + y'(\tau_i)^2} \Delta t_i- \sum_i^n\sqrt{x'(\tau_i)^2 + y'(\tau_i')^2}\Delta t_i \]

利用三角形中的不等式,我们得到
\[\leq \sum_i^n|y'(\tau_i') - y'(\tau_i)|\Delta t_i\] 
由于$x(t),y(t) \in C^1[\alpha,\beta]$
$x'(t),y'(t)$在$[\alpha,\beta]$上一致连续,
\[\forall \varepsilon > 0, |P| < \delta , |\tau_i - \tau'_i| <\delta\]
\[\abs{y'(\tau_i) - y'(\tau_i')} < \dfrac{\varepsilon}{2(\beta - \alpha)}\]

因而前面的不等式继续化简为
\[
\begin{aligned}  
&< \sum_i^n \dfrac{\varepsilon}{2(\beta - \alpha)} \Delta t_i\\
&=\dfrac{\varepsilon}{2(\beta - \alpha)} \sum_i^n \Delta t_i\\
&=\dfrac{\varepsilon}{2(\beta - \alpha)} (\beta - \alpha)\\
&=\dfrac{\varepsilon}{2}\\
\end{aligned}
\]

又由于\[|P| < \delta',\abs{\sum_i^n\sqrt{x'(\tau_i)^2 + y'(\tau_i)^2} - L }< \dfrac{\varepsilon}{2}\]

取$|P| = \min\{\delta,\delta'\}$
我们可以利用绝对值三角不等式证明
\[\abs{\sum_i^n\sqrt{x'(\tau_i)^2 + y'(\tau_i')^2} - L}< \dfrac{\varepsilon}{2} +\dfrac{\varepsilon}{2} = \varepsilon \]
由此可知原式可积且结果为
\[L =\int_{\alpha}^{\beta} \sqrt{x'(t)^2 + y'(t)^2}\]

同理对于极坐标的情况.我们的长度同理可以证明即为黎曼和的极限
\[ L = \int_{\alpha}^{\beta} \sqrt{x(\theta)^2  + y(\theta)^2}\, d\theta\]

\subsubsection{旋转体的侧面积}
同理进行处理,利用函数的一致连续性证明原先的表达式和黎曼和的差距为无穷小

\[\lim_{|P|\to 0} = \sum_{i = 1}^{n} 2\pi y(\xi_i)\sqrt{x'(\xi_i')^2 + y'(\xi_i'')^2} \Delta t = 2 \pi \int_{\alpha}^{\beta}y(t)\sqrt{x'(t)^2 + y'(t)^2}\, dt\]

\subsubsection{线密度$\rightarrow$质量}
\[\Delta m_i = \rho(\xi_i)\sqrt{x'(\xi_i')^2 + y'(\xi_i'')^2}\Delta t_i = \int_{\alpha}^{\beta} \rho(t) \sqrt{x'(t)^2 + y'(t)^2}\, dt\]

\subsection{可积性}
可积:黎曼和极限$\lim_{|P|\to 0} \sigma(f,P,\xi)$存在\\

我们要探究这样的积分存在的条件

首先利用夹逼原理:

$f(x)$在$[a,b]$上有界,作分割$a = x_0 < \cdots < x_n = b$
那么$f(x)$在$[x_{k - 1},x_k]$上有界,而且必有上下确界,

我们记这个区间上的下确界为$m_k = \inf_{x\in[x_{k - 1},x_k]}f(x)$,上确界为$M_k = \sup_{x\in[x_{k - 1},x_k]}f(x)$,$\omega_k = M_k - m_k$称为振幅
,

记$m$为$f$在整个区间上的最小值,$M$为$f$在整个区间上的最大值,$\omega$为f在整个区间上的振幅

\begin{identification}
\[L(f,p) = \sum_{i}^{n}m_i \Delta x_i\]
称为Darboux下和
\[U(f,p) = \sum_{i}^{n}M_i \Delta x_i\]
称为Darboux上和,二者统称为Darboux和
\end{identification}
显然我们有
\[L(f,p) \leq \sigma(f,p,\xi) \leq U(f,p)\]

我们只需要证明Darboux上和等于Darboux下和即可利用夹逼原理\\

下面我们证明这与划分方式无关,只与划分区间的宽度有关

\begin{lemma}
  $P:[a,b], a = x_0 < x_1 \cdots < x_n = b$
  我们增加一个分割点$x'\in[x_{k - 1},x_{k}]$,形成一个新的分割$P'$
于是就有
  \[U(f,P) \geq U(f,P')\]
  \[L(f,P) \leq L(f,P')\]
  加入一个点后
  \[L(f,P') - L(f,P) = (m_k' - m_k)(x' - x_{k - 1}) + (m_k'' - m_k)(x_{k} - x') \leq \omega_k(x_k - x_{k - 1}) \leq \omega_k |P|\]

  因此得出:如果$P'$为$P$增加$l$个分割点得到的分割,那么\[L(f,P) \leq L(f,P') \leq L(f,P) + l ~\omega~|P|\]
  同理\[U(f,P)  - l~ \omega~|P|\leq U(f,P') \leq U(f,P)\]
\end{lemma}

\begin{lemma}
  $P_1,P_2$是$[a,b]$的任意两个分割,一定有$L(f,P_1)\leq U(f,P_2)$\\
\end{lemma}

\begin{cproof}
  合并两个分割记为$P'$,由前面的引理可知
  \[L(f,P_1)\leq L(f,P') \leq U(f,P')\leq U(f,P_2)\]
\end{cproof}

$L(f,P)$有上界,有上确界;同样,$U(f,P)$有下界,有下确界

\[\sup L(f,P) = \underline{I} := \int_{a}^{b}f(x)\, dx\]称为上积分\\
\[\inf U(f,P) = \bar{I} := \int_{a}^{b}f(x)\, dx\]称为下积分\\
联系上下极限一定存在,上下积分也是一定存在的\\

\begin{lemma}
  (1)$\lim_{|P|\to 0}L(f,P) = \underline{I}~$
  (2)$\lim_{|P|\to 0}U(f,P) = \bar{I}$

\end{lemma}

\begin{cproof}
  \[\forall \varepsilon > 0 ,\exists P_0 : a = x_0 < \cdots < x_n = b, ~s.t.~ \underline{I} - \dfrac{\varepsilon}{2} < L(f 
  ,P_0) < \underline{I}\]
  (这一条是根据上确界的定义得出的)

若任意分割$P,|P| < \delta = \dfrac{\varepsilon}{2~l~w + 1}$,
我们把$P_0$和$P$合并得到$P'$,根据引理
\[\underline{I} - \dfrac{\varepsilon}{2} < L(f,P_0) \leq L(f,P')<  \underline{I}\]
而且$P'$为$P$增加了$l$个分点,那么还可以根据引理得到不等式
\[L(f,P) \leq L(f,P') \leq L(f,P) + lw|P|\]
移项得到\[L(f,P) \geq L(f,P') - l w |P| \geq \underline{I} - \dfrac{\varepsilon}{2} - \dfrac{\varepsilon l \omega}{2 l \omega + 1} > \underline{I} - \varepsilon\]
\end{cproof}

\begin{theorem}
  $f(x)$在$[a,b]$上有界,以下三个说法等价

  (1)~$\forall \varepsilon > 0,~\exists |P|,~s.t.~ U(f,P) - L(f,P) < \varepsilon$

  (2)~$\underline{I} = \bar{I}$

  (3)~$f(x)$在$[a,b]$上可积
\end{theorem}

\begin{cproof}
(1)$\Rightarrow$(2)\\
如果$\forall \varepsilon > 0,~\exists |P|,~s.t.~ U(f,P) - L(f,P) < \varepsilon$,那么
\[0 \leq \bar{I} - \underline{I} \leq \varepsilon\]
由$\varepsilon$的任意性可知
\[\bar{I} = \underline{I}\]

(2)$\Rightarrow$(3)\\
$\bar{I} = \underline{I}$,由夹逼原理
\[L(f,P) \leq \sigma(f,P,\xi) \leq U(f,P)\]
可知
\[\lim_{|P| \to 0} \sigma(f,P,\xi) = I\]

(3)$\Rightarrow$(1)\\
\[\lim_{|P| \to 0} \sigma(f,P,\xi) = I\]
用对应语言可以表述为
\[\forall \varepsilon > 0 ,\exists \delta > 0, \forall |P| < \delta , |\sigma(f,P,\xi) - I | \leq \dfrac{\varepsilon}{3}\]
我们对于每个区间的$\xi$分别取到函数值的上下确界
再令两式相减,得到
\[ 0 < U(f,P) - L(f,P) \leq \dfrac{2\varepsilon}{3} < \varepsilon\]
\end{cproof}

\paragraph{等价形式}

(1)~$\forall \varepsilon > 0,\exists P ~s.t. \Omega(f,P) < \varepsilon$

(2)~$\lim_{|P|\to 0} \Omega(f,P) = 0$

(3)~$f\text{在}[a,b]\text{可积}$\\

\begin{example}
  \begin{equation*}
    D(x) = 
    \begin{cases}
      1,x\in \RR\setminus \mathbb{Q}\\
      0,x\in\mathbb{Q}\\
    \end{cases}
  \end{equation*}
对\[\forall x \in [a,b],\Omega(D(x),P) = \sum_{i = 1}^{n} \omega_i \Delta x_i = \sum_{i = 1}^{n}\Delta x_i = b - a\]
\end{example}

\subsection{可积函数类}
\begin{identification}
若$f$在$[a,b]$上黎曼可积,则记
\[f\in R[a,b]\]
$R[a,b]$是一个线性空间,(前面我们知道连续函数空间$C[a,b]$也是一个线性空间)\\
\end{identification}

\begin{theorem}
  $[a,b]$是$[\tilde{a},\tilde{b}]$的任意闭子区间,如果$f\in R[\tilde{a},\tilde{b}]$,则
  \[f\in R[a,b]\]
\end{theorem}

\begin{cproof}
$P$是$[a,b]$的分割,$a = x_0 < x_1 < \cdots < x_n = b$\\

$\forall \varepsilon > 0,P$增加分割点,构成$[\tilde{a},\tilde{b}]$的分割$\tilde{P}$,已知$f\in R[\tilde{a},\tilde{b}]$

\[\exists \delta > 0,|P| < \delta ,|\tilde{P}| < \delta\]

\[\Omega(f,\tilde{P}) < \varepsilon\]
\[\Omega(f,P) \leq\Omega(f,\tilde{P}) < \varepsilon\]

(后者有更多非负的段)
\end{cproof}

\begin{lemma}
  函数$\varphi$在$J$上有定义
  \[\omega(\varphi) = M(\varphi) - m(\varphi)\]
  则
  \[\sup_{x',x''\in J}|\varphi(x') - \varphi(x'')| = \omega(\varphi)\]
\end{lemma}

\begin{cproof}
  第一种情况即 $M(\varphi) = m(\varphi)$证明是显然的故省略

  如果$\omega(\varphi) > 0$,那么我们有
  \[\varphi(x') - \varphi(x'') \leq M(\varphi) - m(\varphi)\]
  \[ \varphi(x'')-\varphi(x')  \leq M(\varphi) - m(\varphi)\]

  这样我们可知道$\omega(\varphi)$是其上界,下面要利用上确界的定义证明其是上确界

  \[\forall \varepsilon > 0 (\varepsilon < \omega(\varphi))\]
  由于$M(\varphi)$是上确界,\[\exists x'\in J ,s.t.\varphi(x') > M(\varphi) - \dfrac{\varepsilon}{2}\]
  由于$m(\varphi)$是下确界,\[\exists x''\in J ,s.t.\varphi(x'') < m(\varphi) + \dfrac{\varepsilon}{2}\]
  因此我们将两式相减即可得到
  \[|\varphi(x') - \varphi(x'')| > \omega(\varphi) - \varepsilon\]
  从而我们证明了上确界
\end{cproof}

\begin{theorem}
  如果  \[f,g\in R[a,b] \Rightarrow fg \in R[a,b]\]
\end{theorem}

\begin{cproof}
  \[\forall x',x''\in J_i = [x_{i - 1},x_i]\]
  \[
  \begin{aligned}  
  |f(x')g(x') - f(x'')g(x'')|
  &=|(f(x')-f(x''))g(x') + f(x'')(g(x') - g(x''))|\\
  &\leq L|(f(x')-f(x''))| + K|g(x') - g(x'')|\\
  \end{aligned}
  \]
  即我们得到
  \[\omega_i(fg) \leq L\omega_i(f) + K\omega_i(g)    \]

  累加得到
  \[\Omega(fg,P) \leq L\Omega(f,P) + K\Omega(g,P)\]

  之后我们取极限$|P|\to 0$即可得到
  \[\lim_{|P|\to 0} \Omega(fg, P) = 0\]
  这说明了$fg$是黎曼可积的
\end{cproof}

\begin{theorem}
  \[f\in R[a,b] \text{且} f(x) > d > 0\Rightarrow \dfrac{1}{f} \in R[a,b]\]
\end{theorem}

\begin{cproof}
  \[
  \begin{aligned}
  \omega_i(\dfrac{1}{f}) &= \sup_{x',x''\in J}\abs{\dfrac{1}{f(x')} - \dfrac{1}{f(x'')}} \\
  &= \sup_{x',x''\in J}\abs{\dfrac{f(x') - f(x'')}{f(x')f(x'')}}\\ 
  &\leq \dfrac{1}{d^2}\sup_{x',x''\in J}|f(x') - f(x'')|\\
  &=\dfrac{1}{d^2}\omega_i(f)    \\
  \end{aligned}
  \]

  累加得到\[\Omega(\dfrac{1}{f},P) \leq  \dfrac{1}{d^2}\Omega(f,P)\]
  取极限$|P|\to 0$即可得到
\end{cproof}

\begin{theorem}
  \[f\in R[a,b] \Rightarrow |f| \in R[a,b]\]
  注意反过来是推不出来的
\end{theorem}

\begin{cproof}
  $\omega_i(|f|) = \sup_{x',x''\in J}\left||f(x')|- |f(x'')|\right| \leq |f(x') - f(x'')| = \omega_i(f)$
  之后累加求极限即可
\end{cproof}

\begin{theorem}
  $f$在$[a,b]$上单调,则$f \in R[a,b]$\\
\end{theorem}

\begin{cproof}
我们考察单调上升的情形,另一种情况略\\

$\forall \varepsilon > 0 $,取$\delta  = \dfrac{\varepsilon}{f(b) - f(a) + 1}$,分割$P$满足$|P| < \delta$
\[
\begin{aligned}
\Omega(f,P) &= \sum_{i = 1}^{n}\omega_i(f)\Delta x_i\\
&\leq \delta \sum_{i = 1}^{n}\omega_i(f) \\
&=  \delta \sum_{i = 1}^{n}(f(x_i) - f(x_{i - 1}))\\
&=  \delta (f(b) - f(a))\\
&\leq \varepsilon\\
\end{aligned}
\]
\end{cproof}

\begin{identification}
有限变差函数:可以表示为两个单调函数之差,因此是黎曼可积的\\
\end{identification}

\begin{theorem}
  \[f(x) \in C[a,b] \Rightarrow f(x) \in R[a,b]\]
\end{theorem}

\begin{cproof}
  因此$f(x)$在$[a,b]$上一致连续,那么就有
  \[\forall \varepsilon > 0,\exists \delta > 0 ,\forall x',x'' \in [a,b],|x' - x'' | < \delta , \abs{f(x') - f(x'')}  < \dfrac{\varepsilon}{b - a}\]
  $P$为$[a,b]$任何一个分割,$J_i = [x_{i - 1} , x_i],|P| < \delta$

\[\omega_i(f) = \sup_{x',x''\in J}|f(x') - f(x'')|\leq \dfrac{\varepsilon}{b - a}\]
  
  \[\Omega(f,P) = \sum_{i = 1}^{n} \omega_i(f) \Delta x_i
  \leq\dfrac{\varepsilon}{b - a} \sum_{i = 1}^{n} \Delta x_i
  = \varepsilon\]
\end{cproof}

\begin{theorem}
  $f$在$[a,b]$上有界,但在有限个点$c_1,\cdots,c_l$外连续,则$f\in R[a,b]$
\end{theorem}

\begin{cproof}
  设间断点$c_1,\cdots ,c_l ,$我们如下构造$c_k\in J_k = (c_k - \eta , c_k + \eta ) $, 且有 $\bigcup_{i = 1}^l J_k $包含所有$c_k$\\

  因而$[a,b]\setminus \bigcup_{i = 1}^l J_k$是一个闭集
  $f(x)$在其上一致连续
  \[\forall \varepsilon > 0,\exists \delta > 0 ,\forall x',x'' \in [a,b]\setminus \bigcup_{i = 1}^l J_k,|x' - x'' | < \delta , \abs{f(x') - f(x'')}  < \eta\]
  作$[a,b]$的分割 $|P| < \min\{\delta,\eta\}$
  与$J_i$相交的划分区间最多四个,长度 $< 4 \times l \times \eta$
  振幅之和可以分为两部分,第一部分与$J_k$不相交,第二部分与$J_k$相交
  \[
  \begin{aligned}
  \Omega(f,P) &= \sum_{i}' \omega_i \Delta x_i + \sum_{j}'' \omega_j \Delta x_j\\
  &\leq \eta \sum_{i}'\Delta x_i + 4l\eta \times  \omega_{max}\\
  &\leq \eta(b - a) + 4 l\eta  \times \omega_{max}\\
  \end{aligned}
  \]
  那么我们只要取 \[0 < \eta < \dfrac{\varepsilon}{b - a + 4 l\eta\times\omega_{max}}\]
  即可证明
\end{cproof}

\begin{theorem}
  $f,g$在$[a,b]$上定义,除了有限点$c_1,\cdots ,c_l$外,$f(x) \equiv g(x),$那么\[f\in R[a,b] \Rightarrow g \in R[a,b] ,\text{且}\int_{a}^{b}f(x)\, dx = \int_{a}^{b}g(x)\, dx \]
\end{theorem}

\begin{cproof}
  分割$P$最多有两个区间包含$c_k$
  \[\abs{\sigma(g,P,\xi) - \sigma(g,P,\xi)} \leq 2lk |P| \rightarrow 0,\text{其中}k = \max_{1\leq k\leq l}\{g(c_k) - f(c_k)\}\]
\end{cproof}

迪利克雷函数改变了有理数上的函数值,而所有有理数的宽度应该是0,它们理应没有面积\\
黎曼积分做的是有限分割,再去考虑极限,如果要真正反映面积,应当要做无穷分割,那就需要无穷覆盖,在实变函数中讨论\\
\subsubsection{复合函数的可积性}
\begin{theorem}
  
  $y = f(x)$在$[a,b]$上黎曼可积,$g(y)$关于$y$可导且导数有界,则$g(f(x))\in R[a,b]$
  
\end{theorem}

\begin{cproof}
\[f\in R[a,b] \Leftrightarrow \sum_{i =1}^{n} \omega_i(f) \Delta x_i \to 0 (|P| \to 0)\]

而对于复合函数
\[\sum_{i =1}^{n} \omega_i g(f(x)) \Delta x_i \]

其中
\[
\begin{aligned}  
  \omega_i(g(f(x))) &= \sup_{x',x''\in[x_{i - 1},x_i]} \abs{g\parameter{f(x')} - g\parameter{f(x'')}}\\
  &= |g'(\eta_i)| \sup_{x',x''\in[x_{i - 1},x_i]} \abs{f(x') - f(x'')}\text{(拉格朗日中值定理)}\\
  &= |g'(\eta_i)| \omega_i(f)\\
\end{aligned}  
  \]


若$|g'(\eta_i)|$有界,那么 \[\sum_{i =1}^{n} \omega_i g(f(x)) \Delta x_i  \to 0(|P| \to 0)\]
从而$g(f(x))\in R[a,b]$
\end{cproof}

\subsection{变上限的积分}

\begin{lemma}
  $f\in R[a,b] ,|f(x)| \leq K ,\forall x\in [a,b]$
  则\[\abs{\int_{a}^{b}f(x)\,dx} \leq K|b - a|\]
\end{lemma}

\begin{cproof}
如果 $a < b, -K \leq f(x) \leq K$\\
\[-K(b - a) \leq \int_{a}^{b}f(x) \, dx \leq K(b - a)\]
\[\abs{\int_{a}^{b}f(x)\,dx} \leq K|b - a|\]
另一种情况同理
\end{cproof}
\begin{theorem}

设$f\in R[a,b],\forall x\in [a,b]$则$ [a,x] \subset[a,b]$构造
\[\varPhi(x) = \int_{a}^{x}f(t)\,dt\],则$\varPhi(x)$在$[a,b]$上连续\\
  
\end{theorem}
\begin{cproof}
\[f\in R[a,b]\Rightarrow \abs{f(t)} \leq K,\forall t\in [a,b]\]

\[\forall x_0 \in [a,b], \abs{\varPhi(x) - \varPhi(x_0)} = \abs{\int_{x_0}^{x} f(t)\,dt}
\leq K|x - x_0|\]
取极限$x\to x_0$ ,$\varPhi(x) \to \varPhi(x_0)$
从而我们证明了这个函数的连续性
\end{cproof}
连续性不需要加条件,但下面我们要讲的可导性需要添加其他条件

\begin{theorem}
设$f\in R[a,b],$ $x_0\in [a,b],f$在$x_0$处连续,则$\varPhi(x)$在$x_0$处可导,且有
\[\varPhi'(x_0) = f(x_0)\]
\end{theorem}

\begin{cproof}
  $f$在$x_0$处连续,$\forall \varepsilon > 0,|t - x_0| <\delta$,有
  \[|f(t) - f(x_0)|< \varepsilon\]
  即有
  \[
  \begin{aligned}  
  \abs{\dfrac{\varPhi(x) - \varPhi(x_0)}{x - x_0} - f(x_0)} &= \abs{\dfrac{1}{x - x_0} \int_{x_0}^{x}f(t)\,dt - f(x_0)}\\
  &= \abs{\dfrac{1}{x - x_0} \int_{x_0}^{x}f(t)- f(x_0)\,dt}\\ 
  &\leq \dfrac{1}{x - x_0} \int_{x_0}^{x}\abs{f(t)- f(x_0)}\,dt\\
  &\leq \varepsilon\\
  \end{aligned}
  \]
\end{cproof}

同理,如果$f$在$[a,b]$整个区间上连续,那么$\varPhi$在$[a,b]$上每一点都可导,且有
\[\varPhi'(x) = f(x),\forall x\in[a,b]\]

因此\textbf{连续函数一定有原函数}(只是不一定能用初等函数形式表示出来,例如)\\
\[e^{-x^2},\quad\dfrac{\sin x}{x},\quad,\sin (x^2)\]
但一定能表示为
\[\int_{a}^{x}f(t)\, dt + C\]
的形式

由此也可以得到N-L公式(最初推导这个公式时就是采用这样变上限积分的思路)

牛顿可积:只需要原函数存在(如果$f$在$[a,b]$上存在原函数$F(x)$,则称其牛顿可积)
\begin{example}
  黎曼可积非牛顿可积
\end{example}
\begin{cproof}
  \begin{equation*}
    f(x) = \begin{cases}
      0,\quad x \neq 0\\
      1,\quad x = 0\\
    \end{cases}
  \end{equation*}
这个函数在$[-1,1]$上黎曼可积,然而根据Darboux定理,其不能有第一类间断点,矛盾!可知其不是一个函数的导数,故其不是牛顿可积的
\end{cproof}
\begin{example}
牛顿可积非黎曼可积
\end{example}
\begin{cproof}
  \begin{equation*}
    F(x) = \begin{cases}
      x^{1 + \alpha}\sin \dfrac{1}{x},\quad x \neq 0(0 < \alpha < 1)\\
      0,\quad x = 0\\
    \end{cases}
  \end{equation*}
  \begin{equation*}
    F'(x) = \begin{cases}
      (1 + \alpha)x^{\alpha}\sin \dfrac{1}{x} -\cos \dfrac{1}{x} \dfrac{1}{x^{1 - \alpha}},\quad x \neq 0(0 < \alpha < 1)\\
      0,\quad x = 0\\
    \end{cases}
  \end{equation*}  
我们考察$f(x) = F'(x)$,其有原函数,因此是牛顿可积的,但这个函数是无界函数,因此其不是黎曼可积的
\end{cproof}

N-L公式左侧要求函数黎曼可积的,右侧要求有原函数(牛顿可积)\\

而前面我们证明了连续函数都有原函数,因此对于连续函数N-L公式一定是成立的(然而,注意可积并不需要有原函数,N-L也不成立)\\
\begin{theorem}
  若$f$在$x_0$处左连续,我们就限制$x < x_0$, $\Rightarrow \varPhi(x)$在$x_0$处左侧可导,
  \[\varPhi'_-(x_0) = f(x_0)\]
  若$f$在$x_0$处右连续,我们就限制$x > x_0$, $\Rightarrow \varPhi(x)$在$x_0$处右侧可导,
  \[\varPhi'_+(x_0) = f(x_0)\]
  如果对于$[a,b](a< b)$,如果$f$在$a$点右连续,那么$\varPhi$在$a$点右导数 \[\varPhi'_+(a) = f(a)\]
  \[\varPhi'_-(b) = f(b)\]
\end{theorem}

\subsubsection{变下限积分}

$\Psi(x) = \int_{x}^{b} f(t)\, dt = -\int_{b}^{x} f(t)\, dt $,
若$f$在$x_0$处连续,则
\[\Psi'(x_0) = -f(x_0)\]

如果$\Psi$在$a$点右连续\[\Rightarrow \Psi'_+(a) = -f(a)\]
如果$\Psi$在$b$点左连续\[\Rightarrow \Psi'_-(b) = -f(b)\]

\subsubsection{复合函数}
$F(x) = \int_{u(x)}^{v(x)} f(t)\, dt $,其中$u(x),v(x)\in[a,b],f\in C[a,b]$,则
\[F(x) = \int_{c}^{v(x)} f(t)\, dt + \int_{u(x)}^{c} f(t)\, dt\]
\[F'(x) = f(v(x)) v'(x) - f(u(x))u'(x)\]

\subsubsection{变上限积分的应用-平面上曲线的曲率和曲率半径}

曲率-曲线的弯曲程度

\begin{equation*}
  \begin{cases}
    x = x(t)\\
    y = y(t)\\
  \end{cases}
\end{equation*}

从$M$到$M_1$,曲线切线方向的变化$\alpha \to \alpha_1,\Delta \alpha = \alpha_1 - \alpha$
曲率:切线角度相对于弧长的改变率

\[
\begin{aligned}
\lim_{\Delta s \to 0}\abs{\dfrac{\Delta \alpha}{\Delta s}}
&= \lim_{\Delta s \to 0}\abs{\dfrac{\frac{\Delta \alpha}{\Delta t}}{\frac{\Delta s}{\Delta t}}}\\
&= \abs{\dfrac{\frac{d \alpha}{d t}}{\frac{d s}{d t}}}\\
\end{aligned}
\]
而$\tan \alpha = \dfrac{y'(t)}{x'(t)} = \dfrac{dy}{dx}$,即有
\[\alpha = \arctan \dfrac{y'(t)}{x'(t)}\]
即有\[\dfrac{d \alpha}{d t} =\dfrac{x'(t) y''(t) - x''(t) y'(t)}{x'(t)^2 + y'(t)^2}\]

同理利用\[S(t) = \int_{a}^{t} \sqrt{x'(t)^2 + y'(t)^2}\]
可得\[\dfrac{ds}{dt}= \sqrt{x'(t)^2 + y'(t)^2}\]

从而我们得到曲率
\[\abs{\dfrac{x'(t) y''(t) - x''(t) y'(t)}{(x'(t)^2 + y'(t)^2)^{\frac{3}{2}}} }= \kappa\] 
我们将其倒数
\[R = \dfrac{1}{\kappa}\]
定义为曲率半径\\

曲线上有关一阶导和二阶导的信息可使用圆来考虑\\

取 $t = x$,则$\kappa =  \dfrac{\abs{y''(x)}}{(1 + y'(x)^2)^{\frac{3}{2}}}$

\subsection{积分不等式}

\begin{theorem}
  Young不等式

  $\varphi(x)$在$[0,C]$上严格单调且连续,则其反函数$\varPsi(x)$同样是严格单调且连续的,如果$\varphi(0) = 0$,任取$a \in [0,C],b\in[0,\varphi(C)],\varphi\in R[0,a],\varPhi \in R[0,b]$,则
  \[ab \leq \int_{0}^{a} \varphi(x) \, dx + \int_{0}^{b} \varPsi(y)\, dy\]
  (其中当且仅当$b = \varphi(a)$时取到等号)
\end{theorem}
\begin{cproof}
严格证明是可能的(可以参见教材)
\end{cproof}

\begin{lemma}
  
  $f(x) = \ln x , f''(x) = -\dfrac{1}{x^2}$,因此$f(x)$是一个凹函数$\lambda_1 ,\lambda_2 > 0 ,\lambda_1 + \lambda_2 =1$ \\
  \[\forall x,y > 0, \lambda_1 \ln x + \lambda_2 \ln y \leq \ln(\lambda_1 x + \lambda_2 y)\]
  
  进而我们化简可以得到\[x^{\lambda_1} y^{\lambda_2} \leq \lambda_1 x + \lambda_2 y\]
  我们令$\lambda_1 = \dfrac{1}{p},\lambda_2 = \dfrac{1}{q}$,从而不等式化为
  \[x^{\frac{1}{p}} y^{\frac{1}{q}} \leq \dfrac{1}{p}x + \dfrac{1}{q} y,\forall x,y >0 (x = y\text{时取等号})\]
  
\end{lemma}
我们可以继续推广这个不等式,得到离散形式的$H\ddot{o}lder$不等式,

\begin{theorem}
  离散形式的$H\ddot{o}lder$不等式
  \[\sum_{i = 1}^n a_i b_i \leq \parameter{\sum_{i = 1}^n a_i^p}^{\frac{1}{p}} \parameter{\sum_{i = 1}^n b_i^q}^{\frac{1}{q}}\]
\end{theorem}

\begin{cproof}
取\[x_i = \dfrac{a_i^p}{\sum_{i = 1}^n a_i^p} , y_i = \dfrac{b_i^q}{\sum_{i =1 }^n b_i^q}\]
带入可得
\[\dfrac{a_i b_i}{\parameter{\sum_{i =1}^n  a_i^p}^{\frac{1}{p}}\parameter{\sum_{i = 1}^n b_i^q}^{\frac{1}{q}}}\leq \dfrac{1}{p}\dfrac{a_i^p}{\sum_{i = 1}^n a_i^p} + \dfrac{1}{q} \dfrac{b_i^q}{\sum_{i = 1}^n b_i^q}\]

将这个式子对$i$累加,并移项可以得到
\[\sum_{i = 1}^n a_i b_i \leq \parameter{\sum_{i = 1}^n a_i^p}^{\frac{1}{p}} \parameter{\sum_{i = 1}^n b_i^q}^{\frac{1}{q}}\]

\end{cproof}

事实上,我们可以从这个式子直接推广到积分形式

取$a_i = |f(\xi_i)|\Delta x_i^{\frac{1}{p}} , b_i = |g(\xi_i)|\Delta x_i ^\frac{1}{q}$即可,我们得到黎曼和,之后转换为积分形式即可\\

\begin{theorem}

积分形式的$H\ddot{o}lder$不等式

$p > 1,q > 1 , \dfrac{1}{p } + \dfrac{1}{q} = 1,f \in R[a,b],g\in R[a,b]$,则
\[\abs{\int_{a}^{b}fg\, dx} \leq \parameter{\int_{a}^{b}|f^p(x)|\,dx}^{\frac{1}{p}} \parameter{\int_{a}^{b}|g^q(x)|\,dx}^{\frac{1}{q}}\]
\end{theorem}

\begin{cproof}
不妨设$\int_{a}^{b}|f^p(x)|\,dx > 0, \int_{a}^{b}|g^q(x)|\,dx > 0$,取
\[\varphi(x) = \dfrac{f(x)}{\parameter{\int_{a}^{b}|f^p(x)|\,dx}^{\frac{1}{p}}},\varPsi(x) = \dfrac{g(x)}{\parameter{\int_{a}^{b}|g^q(x)|\,dx}^{\frac{1}{q}}}\]
\[(\text{注意此时的分母其实相当于为常数})\]
之后在不等式$x^{\frac{1}{p}} y^{\frac{1}{q}} \leq \dfrac{1}{p}x + \dfrac{1}{q} y$中令$x = |\varphi(x)|^p, y = |\varPsi(x)|^q$,得到
\[|\varphi(x)\varPsi(x)| \leq \dfrac{|\varphi(x)|^p}{p} + \dfrac{|\varPsi(x)|^q}{q}\]
之后我们对两边积分
\[\int_{a}^{b}|\varphi(x)\varPsi(x)|\, dx \leq \int_{a}^{b}\dfrac{|\varphi(x)|^p}{p}\,dx +\int_{a}^{b} \dfrac{|\varPsi(x)|^q}{q}\, dx\]

之后我们再带入即可,注意到右式实际上就是$\dfrac{1}{p } + \dfrac{1}{q} = 1$,即\\
\[ \dfrac{\abs{\int_{a}^{b}fg\, dx}}{\parameter{\int_{a}^{b}|f^p(x)|\,dx}^{\frac{1}{p}}\parameter{\int_{a}^{b}|g^q(x)|\,dx}^{\frac{1}{q}}}\leq 1\]
之后我们再进行移项即可
\end{cproof}

同样我们可以把积分形式转化为离散形式,只需要设$f,g$为阶梯函数,$f(x) = a_k , x\in \left[\dfrac{k - 1}{n}, \dfrac{k}{n}\right),g(x) = b_k , x\in \left[\dfrac{k - 1}{n}, \dfrac{k}{n}\right)$,我们取划分点$x_k = \dfrac{k}{n}$同时取标志点$\xi_k = \dfrac{k}{n}$\\

\begin{theorem}
  Minkowski不等式
  \[\parameter{\int_{a}^{b} \abs{f(x) + g(x)}^p\, dx}^{\frac{1}{p}}  \leq \parameter{\int_{a}^{b}|f(x)|^p\,dx}^{\frac{1}{p}}+\parameter{\int_{a}^{b}|g(x)|^p\,dx}^{\frac{1}{p}}\]
\end{theorem}

这个不等式本质上是三角不等式的推广

\begin{cproof}

$p = 1$时即为三角不等式,证明略\\

$p > 1$时,

我们考察两个函数$\varphi(x) = |f(x)|$和$\varPsi(x) = \abs{f(x) + g(x)}^{\frac{p}{q}}$,利用$H\ddot{o}lder$不等式
\begin{equation}
  \int_{a}^{b}\abs{f(x)}\abs{f(x)+ g(x)}^{\frac{p}{q}} \leq \parameter{\int_{a}^{b}\abs{f(x)}^p\, dx}^{\frac{1}{p}} \parameter{\int_{a}^{b}\abs{f(x) + g(x)}^p}^{\frac{1}{q}}
\end{equation}
  同理对于$g(x)$可以得到同样的不等式
\begin{equation}
  \int_{a}^{b}\abs{g(x)}\abs{f(x)+ g(x)}^{\frac{p}{q}} \leq \parameter{\int_{a}^{b}\abs{g(x)}^p\, dx}^{\frac{1}{p}} \parameter{\int_{a}^{b}\abs{f(x) + g(x)}^p}^{\frac{1}{q}}
\end{equation}

利用\[1 + \dfrac{p}{q} = p\]
我们可以得到
\[
\begin{aligned}
  \parameter{\int_{a}^{b} \abs{f(x) + g(x)}^p\, dx} &= \int_{a}^{b}\abs{f(x) + g(x)}^{1 + \frac{p}{q}}\,dx \\
  &= \int_{a}^{b}\abs{f(x) + g(x)}\abs{f(x)+ g(x)}^{\frac{p}{q}} \\
  &\leq \abs{\parameter{\int_{a}^{b}\abs{f(x)}^p\, dx}^{\frac{1}{p}} + \parameter{\int_{a}^{b}\abs{g(x)}^p\, dx}^{\frac{1}{p}} }\times \parameter{\int_{a}^{b}\abs{f(x) + g(x)}^p}^{\frac{1}{q}}\\
\end{aligned}  
  \]
  这里把等式右边最后一项除过来即可
  

\end{cproof}

\begin{theorem}
  离散形式的Minkowski不等式
  \[\parameter{\sum_{i = 1}^{n} \abs{x_k + y_k}^p}^{\frac{1}{p}} \leq \parameter{\abs{x_k}^p}^{\frac{1}{p}} + \parameter{\abs{y_k}^p}^{\frac{1}{p}}\]
\end{theorem}

p = 1即三角不等式

p = 2即柯西不等式

\begin{cproof}
证明思路大体相同,
\[\sum_{i = 1}^{n}\abs{x_k}\abs{x_k + y_k}^{\frac{p}{q}} \leq \parameter{\sum_{i = 1}^{n} \abs{x_k}^p}^{\frac{1}{p}}
\parameter{\sum_{i = 1}^{n} \abs{x_k + y_k}^p}^{\frac{1}{p}}\]

同理仿照上式可以得到

\[\sum_{i = 1}^{n} \abs{y_k}\abs{x_k + y_k}^{\frac{p}{q}} \leq \parameter{\sum_{i = 1}^{n} \abs{y_k}^p}^{\frac{1}{p}}
\parameter{\sum_{i = 1}^{n} \abs{x_k + y_k}^p}^{\frac{1}{p}}\]


\[
\begin{aligned}
  \sum_{i = 1}^{n}\abs{x_k + y_k}^p &= \sum_{i = 1}^{n}\abs{x_k + y_k}^{1 + \frac{p}{q}} \\
  &\leq \sum_{i = 1}^{n} \abs{x_k}\abs{x_k + y_k}^{\frac{p}{q}}+\sum_{i = 1}^{n}\abs{y_k}\abs{x_k + y_k}^{\frac{p}{q}}\\
  &\leq \bracket{\parameter{\sum_{i = 1}^{n} \abs{x_k}^p}^{\frac{1}{p}} + \parameter{\sum_{i = 1}^{n} \abs{y_k}^p}^{\frac{1}{p}}} \parameter{\sum_{i = 1}^{n} \abs{x_k + y_k}^p}^{\frac{1}{p}}\\
\end{aligned}  
  \]

\end{cproof}

同样我们可以实现对于连续形式到离散形式的转化,我们只要按照前面的方式构造一个阶梯函数即可

\subsubsection{凹凸性与积分不等式}

\begin{lemma}
  离散形式的Jensen不等式

  $f(x)$在$[\alpha,\beta]$上是凸函数,$x_i \in[\alpha,\beta]$
  对于权重之和$\lambda_i > 0 \sum_{i = 1}^{n}\lambda_i = 1$

  \[f\parameter{\lambda_1 x_1 + \cdots + \lambda_n x_n} \leq\sum_{i = 1}^{n}\lambda_i f(x_i)\]
\end{lemma}

我们利用其推导连续形式
\begin{theorem}
  连续形式的Jensen不等式

  $f(x)$在$[\alpha,\beta]$上是凸函数.设$\varphi(x) \in R[a,b]$,其值域包含于$[\alpha,\beta]$,$p(x)\in R[a,b],p(x) > 0$,则
  \[f\parameter{\dfrac{\int_{\alpha}^{\beta} p(x) \varphi(x)\,dx}{\int_{\alpha}^{\beta} p(x) \, dx}} \leq \dfrac{\int_{\alpha}^{\beta} p(x)f(\varphi(x)) \, dx }{\int_{\alpha}^{\beta}p(x) \,dx}\]
\end{theorem}
\begin{cproof}
\[\forall p_i > 0, i = 1, \cdots ,n\]
取\[\lambda_i = \dfrac{p_i}{\sum_{i =1 }^{n} p_i},\text{故而有}\sum_{i = 1}^{n}\lambda_i = 1,\]

代入离散形式的Jensen不等式,得到
\[f\parameter{\dfrac{\sum_{i =1 }^{n} p_i x_i}{\sum_{i =1 }^{n} p_i}} \leq \dfrac{\sum_{i =1 }^{n} p_i f(x_i)}{\sum_{i =1 }^{n} p_i}\]
其中$x_i \in [\alpha,\beta]$\\

取$[a,b]$的分割点$a = t_0 < t_1 < \cdots < t_n = b$

\[x_i = \varphi(t_i), p_i = p(t_i) \Delta t_i\]
代入
\[f\parameter{\dfrac{\sum_{i =1 }^{n} p(t_i)  \varphi(t_i) \Delta t_i}{\sum_{i =1 }^{n} p(t_i) \Delta t_i}} \leq \dfrac{\sum_{i =1 }^{n} p(t_i)  f(\varphi(t_i)) \Delta t_i}{\sum_{i =1 }^{n} p(t_i) \Delta t_i}\]

由于$f$是凸函数,其是连续的(极限符号可以拿到$f$里面去),那么令$|P| \to 0$,得到
\[f\parameter{\dfrac{\int_{\alpha}^{\beta} p(x) \varphi(x)\,dx}{\int_{\alpha}^{\beta} p(x) \, dx}} \leq \dfrac{\int_{\alpha}^{\beta} p(x)f(\varphi(x)) \, dx }{\int_{\alpha}^{\beta}p(x) \,dx}\]

如果是凹函数反号即可
\end{cproof}

这个公式的主要记忆方式是将$p(x)$理解为某种加权方式,只是加权不再是简单的求和而是以积分形式,就像我们定义多项式空间中的内积时,可以选择离散的一系列点,也可以选择区间上的积分

\begin{theorem}
  \[\lim_{p \to + \infty} \parameter{\int_{a}^{b} |f(x)|^p \, dx}^{\frac{1}{p}} = \max_{x\in[a,b]}\abs{f(x)}\]
\end{theorem}

\begin{cproof}
利用
\[\lim_{n \to +\infty} \parameter{\dfrac{x_1^n + \cdots + x_n^n}{n}}^{\frac{1}{n}} = x_{max}\]
作分割$P$,宽度$\Delta x_i = \dfrac{b - a}{n}$,

取$X_i = M_i = \max_{x\in[x_{i - 1}, x_i]}\abs{f(x)}$\\
代入原等式
\[
\begin{aligned}
  \lim_{p \to +\infty} \parameter{\sum_{i = 1}^{n} M_i^p \Delta x_i}^{\frac{1}{p}} &= \lim_{p \to +\infty} \parameter{\dfrac{M_1^p +\cdots +M_n^p}{n}}^{\frac{1}{p}} (b - a)^{\frac{1}{p}}\\
  &= \max_{i = 1,\cdots ,n}\{M_i\}\\
  &= \max_{x\in[a,b]}\abs{f(x)}\\
\end{aligned}  
  \]

取$X_i = m_i = \min_{x\in[x_{i - 1}, x_i]}\abs{f(x)}$\\
同理可以证明
\[
\begin{aligned}
  \lim_{p \to +\infty} \parameter{\sum_{i = 1}^{n} M_i^p \Delta x_i}^{\frac{1}{p}} &= \lim_{p \to +\infty} \parameter{\dfrac{m_1^p +\cdots +m_n^p}{n}}^{\frac{1}{p}} (b - a)^{\frac{1}{p}}\\
  &= \max_{i = 1,\cdots ,n}\{m_i\}\\
\end{aligned}  
  \]

  对于一个类似于Darboux上和和下和的形式
  \[\parameter{\sum_{i = 1}^{n} m_i^p \Delta x_i}^{\frac{1}{p}} \leq  \parameter{\int_{a}^{b} |f(x)|^p \, dx}^{\frac{1}{p}} \leq  \parameter{\sum_{i = 1}^{n} M_i^p \Delta x_i}^{\frac{1}{p}}\]
两边取极限

\[\lim_{p \to +\infty} \parameter{\sum_{i = 1}^{n} m_i^p \Delta x_i}^{\frac{1}{p}} \leq \lim_{p \to + \infty} \parameter{\int_{a}^{b} |f(x)|^p \, dx}^{\frac{1}{p}} \leq \lim_{p \to +\infty} \parameter{\sum_{i = 1}^{n} M_i^p \Delta x_i}^{\frac{1}{p}}\]
即
\[\max_{i = 1,\cdots ,n}\{m_i\} \leq \lim_{p \to + \infty} (\int_{a}^{b} |f(x)|^p \, dx)^{\frac{1}{p}} \leq \max_{x\in[a,b]}\abs{f(x)}\]

由于$f$一致连续,故
\[\forall \varepsilon> 0 ,\exists \delta >0, s.t.\forall x,y |x - y| < \delta , |f(x) - f(y)| < \varepsilon \]
\[\text{取}|P| \leq \delta ,M_i - m_i < \varepsilon ,\text{不妨记}i = j\text{时}M_i\text{取到最大值,故}\]
\[
\begin{aligned}\max_{i = 1,\cdots ,n}\{M_i\} - \max_{i = 1,\cdots ,n}\{m_i\} &= M_j - \max_{i = 1,\cdots ,n}\{m_i\}\\
  &\leq M_j - m_j\\
  &< \varepsilon\\
\end{aligned}  
\]
从而

\[\max_{x\in[a,b]}\abs{f(x)} - \varepsilon \leq \lim_{p \to + \infty} \parameter{\int_{a}^{b} |f(x)|^p \, dx}^{\frac{1}{p}} \leq \max_{x\in[a,b]}\abs{f(x)}\]

由于$\varepsilon$的任意性,我们得到
\[\lim_{p \to + \infty} \parameter{\int_{a}^{b} |f(x)|^p \, dx}^{\frac{1}{p}} = \max_{x\in[a,b]}\abs{f(x)}\]
\end{cproof}

\begin{example}
\[\lim_{n \to +\infty} \int_{0}^{2 \pi} f(x)\cos nx \, dx = 0, f(x) \in C[a,b],n\in \NN\]  
\end{example}


\begin{cproof}
构造一个阶梯函数

\begin{equation*}
  H_n(x) =
  \begin{cases}
    f(\dfrac{k -1}{n}  2\pi),\forall x \in \left[\dfrac{k -1}{n}  2\pi,\dfrac{k}{n} 2\pi\right)\\
    f(2 \pi) , x = 2 \pi\\
  \end{cases}
\end{equation*}

这个函数在每个区间上是一个常数$C_k$,而这每个区间恰好横跨$\cos nx$的半个周期,其中一半为正一半为负,因此

\[\int_{0}^{2 \pi} H_n(x) \cos nx \, dx = 0\]

同时利用$f(x)$的一致连续性

\[\forall \varepsilon> 0 ,\exists \delta > 0,\text{对于充分大的}n,\dfrac{2\pi}{n} < \delta\]

\[f(x) - H_n(x) < \dfrac{\varepsilon}{2 \pi}\]

\[\int_{0}^{2 \pi}\bracket{f(x) - H_n(x)}\, dx < |\int_{0}^{2 \pi}\dfrac{\varepsilon}{2 \pi}\,dx| = \varepsilon\]

\[
\begin{aligned}
  \abs{\lim_{n \to +\infty} \int_{0}^{2 \pi} f(x)\cos nx \, dx} &\leq \abs{\lim_{n \to +\infty} \int_{0}^{2 \pi} H_n(x)\cos nx \, dx} + \abs{\lim_{n \to +\infty} \int_{0}^{2 \pi} \bracket{f(x) - H_n(x)}\cos nx \, dx}\\
  &< 0 + \varepsilon\\
  &= \varepsilon
\end{aligned}  
\]

由$\varepsilon$的任意性可知成立
\end{cproof}


\subsubsection{四大连续性}
前两种就是我们常见的连续和一致连续性

\paragraph{Lipschitz连续性}
\[\abs{f(x) - f(y)} \leq L|x - y|,\forall x,y \in [a,b]\]
满足上式成立的最小的$L$称为Lipschitz常数
\begin{theorem}
  若$f\in C^1[a,b]$,则$f$具有Lipschitz连续性
\end{theorem}
\begin{cproof}
$f\in C^1[a,b]$,则$ \abs{f'(\xi)} \leq L$
$\abs{f(x) - f(y)} \leq \abs{f'(\xi)}|x - y|\leq L|x - y|$
\end{cproof}
\paragraph{H$\ddot{o}$lder连续性}

\[\abs{f(x) - f(y)} \leq L|x - y|^{\alpha},\forall x,y \in [a,b],0 < \alpha < 1\]


\begin{theorem}
  $f(x)$和$f'(x)$在$[a,b]$上可积,则$f$具有$H\ddot{o}lder$连续性
\end{theorem}

\begin{cproof}

\[
\begin{aligned}
  \abs{f(x) - f(y)} &= \abs{\int_{y}^{x} f'(t) \times 1 \,dt}&\text{(N-L公式)}\\
  &\leq \bracket{\int_{y}^{x} |f'(x)|^p \, dx}^{\frac{1}{p}} \bracket{\int_{y}^{x} 1^q \, dx}^{\frac{1}{q}}&\text{($H\ddot{o}lder$不等式的积分形式)}\\
  &\leq C \abs{x - y}^{\frac{1}{q}} \\
\end{aligned}  
\]
其中令 $L = C$,~$0 < \alpha = \dfrac{1}{q} < 1$,即可证明定理开始的结论
\end{cproof}

\end{CJK}
\end{document}
\begin{itemize} 
\end{itemize}